\documentclass[11pt,a4paper]{article}

\usepackage{amsmath,amssymb,amsthm,mathtools}
\usepackage[utf8]{inputenc}
\usepackage[margin=1in]{geometry}
\usepackage{hyperref}
\usepackage{enumitem}
\usepackage{cleveref}

\theoremstyle{plain}
\newtheorem{theorem}{Theorem}[section]
\newtheorem{proposition}[theorem]{Proposition}
\newtheorem{lemma}[theorem]{Lemma}
\newtheorem{corollary}[theorem]{Corollary}

\theoremstyle{definition}
\newtheorem{definition}[theorem]{Definition}
\newtheorem{example}[theorem]{Example}
\newtheorem{assumption}[theorem]{Assumption}

\theoremstyle{remark}
\newtheorem{remark}[theorem]{Remark}

\newcommand{\Id}{\mathrm{Id}}
\newcommand{\E}{\mathbb{E}}
\newcommand{\R}{\mathbb{R}}
\newcommand{\PP}{\mathbb{P}}
\newcommand{\QQ}{\mathbb{Q}}
\newcommand{\F}{\mathcal{F}}
\newcommand{\HH}{\mathcal{H}}
\newcommand{\op}{\mathrm{op}}
\newcommand{\im}{\mathrm{im}}
\newcommand{\Ntilde}{\widetilde{N}}
\newcommand{\ip}[2]{\langle #1, #2 \rangle}
\newcommand{\norm}[1]{\| #1 \|}

\crefname{theorem}{Theorem}{Theorems}
\crefname{proposition}{Proposition}{Propositions}
\crefname{lemma}{Lemma}{Lemmas}
\crefname{corollary}{Corollary}{Corollaries}
\crefname{definition}{Definition}{Definitions}
\crefname{assumption}{Assumption}{Assumptions}
\crefname{remark}{Remark}{Remarks}
\crefname{section}{Section}{Sections}

\title{The Boundary of Hedgeability:\\[0.3em]
Pricing and Hedging in Volterra--L\'evy Markets}

\author{Ramiro Fontes\\[0.5em]
\small Quijotic Research\\[0.3em]
\small\texttt{ramirofontes@gmail.com}}

\date{\today}


\begin{document}

\maketitle

\begin{abstract}
We develop the pricing and hedging theory for financial markets
driven by Volterra--L\'evy processes---encompassing rough
volatility models with jumps.  The operator derivative
$D = \delta^*$, defined by duality with the stochastic integral
on a combined Hilbert~$\oplus$~Banach energy space, provides a
single calculus for the continuous and jump components without
Gaussian assumptions.

The central result is a two-parameter Volterra risk-neutral
constraint (VRNC): risk-neutral pricing in markets with memory
and jumps requires solving a Volterra integral equation for an
effective drift rate~$\Psi$ that couples the Brownian and jump
risk premia through the kernel~$K$.  The kernel enters the
pricing structure itself---continuous and jump parameters
cannot be chosen independently.  This holds for all kernel
regularities~$H$, with no semimartingale restriction on the
Volterra process.

The Leibniz defect of the operator derivative---the precise
content invisible to Malliavin calculus---equals the
structurally unhedgeable component of risk.  Its $L^p$~norm
scales as
$T^{2H-1+1/p} \cdot (\int_{|z|\leq 1}|z|^{2p}\,\nu(dz))^{1/p}$:
roughness and jumps interact multiplicatively.  The scaling is
sharp for the variance swap payoff (two-sided BDG bounds).
The roughness parameter~$H$ is proved to be
measure-invariant, establishing it as a structural feature
of the market rather than a pricing artifact.

\medskip\noindent\textbf{MSC 2020:} 60H07, 60H05, 60G51,
60G22, 91G20.\\
\textbf{Keywords:} Volterra--L\'evy processes; operator
derivative; risk-neutral pricing; rough volatility; hedgeability;
Leibniz defect; bracket invariance; calibration invariance.
\end{abstract}


%======================================================================
\section{Introduction}\label{sec:intro}
%======================================================================

\subsection{The pricing problem for rough-jump markets}

What does consistent pricing look like when the market has both
memory and jumps?

Two extensions of classical diffusion models have dominated recent
financial modelling.  \emph{Rough volatility}---volatility driven
by fractional kernels $K_H(t,s) = (t-s)^{H-1/2}$ with
$H < 1/2$---fits the implied volatility surface with remarkable
parsimony~\cite{ElEuchRosenbaum2019,Fukasawa2021}.  Separately,
\emph{L\'evy-driven models} capture the heavy tails and
discontinuities observed in
returns~\cite{ContTankov2004,Applebaum2009}.  But these two
extensions have been studied almost entirely in isolation.  The
Malliavin calculus that underpins rough volatility theory is built
on Wiener space and does not see jumps; the Picard or
Nualart--Vives difference operators for jump processes do not
interact with kernel convolution.  No unified framework exists
for the interaction of roughness and jumps through a
Volterra kernel.

This paper constructs that framework.  The starting point is
the observation that the structurally unhedgeable component of
risk in a Volterra--L\'evy market is \emph{not} the sum of a
rough-volatility error and a jump error.  The two interact
multiplicatively: the kernel amplifies jump risk, and the
amplification is controlled by a single algebraic object---the
kernel-weighted Leibniz defect---that emerges naturally from
the operator factorization program
of~\cite{Fontes2026a,Fontes2026d}.

\subsection{The operator derivative as the right tool}

The operator derivative $D = \delta^*$ is defined by duality
with the stochastic integral~$\delta$, not by any Gaussian
structure~\cite{Fontes2026a}.  For a Volterra--L\'evy
process with both continuous and jump components, $\delta$
acts on the combined energy space $\mathcal{H}_{VL} = H_W
\times \widetilde{H}_J$, and $D = \delta^*$ splits
automatically into a Hilbert derivative~$D_W$ (for the
Brownian part) and a Banach derivative~$\widetilde{D}_J$
(for the jumps).  The key structural fact is that $D$
satisfies the Leibniz rule if and only if the driver has no
jumps~\cite{Fontes2026d}.  In the presence of jumps, the
failure of the Leibniz rule---the \emph{defect}
$\Gamma(F,G) = D(FG) - F \cdot DG - G \cdot DF$---is
precisely the non-Gaussian content.  For Volterra--L\'evy
processes, the kernel $K$ weights this defect, producing
the kernel-weighted Leibniz defect $\Gamma_V^f$ that
controls both the pricing structure and the hedging gap.

The Malliavin derivative on Wiener space and the Picard
difference operator on Poisson space are defined by separate
mechanisms; there is no Malliavin derivative for the combined
process that recovers the kernel-weighted defect as a single
algebraic object.  The operator derivative, defined by duality
with the combined stochastic integral, does this automatically:
the splitting into continuous and jump components is inherited
from the energy space decomposition, and the Leibniz
failure---the non-Gaussian content---emerges as a theorem,
not a definition.


\subsection{Main results}

\begin{enumerate}[label=\textbf{(\Alph*)}]
\item \textbf{The fundamental pricing constraint}
  (\Cref{thm:VRNC-VL}).  In a Volterra--L\'evy market,
  risk-neutral pricing requires solving a Volterra integral
  equation for an effective drift rate
  $\Psi_t = -\sigma_L\,\theta_t^W + \int z(e^{\theta_t^J(z)}-1)\,
  \nu(dz)$
  that aggregates Brownian and jump risk premia through the
  kernel~$K$.  The kernel enters the pricing
  structure itself: continuous and jump pricing parameters
  cannot be chosen independently.  The continuous case
  (no jumps) is recovered as \Cref{cor:VRNC-recovery}.
  The constraint holds for all kernel
  regularities~$H$.

\item \textbf{Structural invariance of the roughness
  parameter}
  (\Cref{thm:bracket-inv}).  The continuous bracket of
  the partial process $V_T(t)$ is deterministic and
  identical under $\mathbb{P}$ and $\mathbb{Q}$:
  the roughness parameter~$H$ estimated from historical
  data is the same~$H$ that governs risk-neutral dynamics.
  This is a structural feature of the energy space, not
  a modelling assumption.

\item \textbf{The hedging gap is the Leibniz defect}
  (\Cref{thm:hedging-main,thm:scaling}).  The
  structurally unhedgeable component of any European
  claim, after delta-hedging with the stock and
  vega-hedging with a variance swap, is the
  kernel-weighted Leibniz defect $\Gamma_V^f$.  Its
  $L^p$~norm scales as
  $T^{2H - 1 + 1/p} \cdot (2p - \gamma)^{-1/p}$:
  roughness and jump intensity interact multiplicatively.
  The scaling is sharp for the variance swap payoff
  (two-sided BDG bounds).  The defect vanishes if and
  only if the driver has no jumps, recovering
  perfect hedging in pure rough volatility
  models~\cite{ElEuchRosenbaum2019}.

\item \textbf{The $\nu$-singularity obstruction persists}
  (\Cref{prop:VT-singular}).  The Volterra process
  $V_T$ is $\nu$-singular:
  $V_T \notin \mathrm{im}(\widetilde{\delta}_J)$.
  The kernel cannot regularize the jump-size singularity.
  This is an intrinsic property of the L\'evy measure,
  independent of the temporal structure.
\end{enumerate}


\subsection{Organization}

\Cref{sec:setup}: the Volterra--L\'evy process and its
semimartingale structure.
\Cref{sec:energy}: the combined energy space and operator
derivatives.
\Cref{sec:ito}: the Volterra--L\'evy It\^o formula and the
kernel-weighted defect decomposition.
\Cref{sec:pricing}: the two-parameter VRNC and bracket
invariance.
\Cref{sec:scaling}: the scaling law for fractional kernels.
\Cref{sec:hedging}: application to hedging in rough-jump models.
\Cref{sec:numerics}: numerical verification.

\section{The Volterra--L\'evy Process}\label{sec:setup}
%======================================================================

\subsection{Probabilistic setup}

Let $(\Omega, \F, (\F_t)_{t \in [0,T]}, \PP)$ be a filtered
probability space supporting:
\begin{enumerate}[label=(\roman*)]
\item a Brownian motion $(W_t)_{t \in [0,T]}$;
\item an independent Poisson random measure $N(dt,dz)$ on
  $[0,T] \times \R$ with intensity
  $dt \otimes \nu(dz)$, where $\nu$ is a L\'evy measure
  satisfying $\int_\R (|z|^2 \wedge 1)\,\nu(dz) < \infty$.
\end{enumerate}
Write $\Ntilde(dt,dz) := N(dt,dz) - dt\,\nu(dz)$ for the
compensated Poisson measure.

\begin{definition}[L\'evy driver]\label{def:levy-driver}
The \emph{L\'evy driver} is
\begin{equation}\label{eq:levy-driver}
  L_t := \sigma W_t + \int_0^t\!\!\int_\R z\,\Ntilde(ds,dz),
  \qquad \sigma \geq 0,
\end{equation}
where the integral converges in $L^p$ for $p \in (1,\gamma)$ when
$\nu$ is the measure of a symmetric $\gamma$-stable process
($\gamma \in (1,2)$), and in $L^2$ when
$\int |z|^2\,\nu(dz) < \infty$.
\end{definition}

\begin{assumption}[Standing assumptions on $\nu$]\label{ass:nu}
Throughout, $\nu$ is symmetric ($\nu(A) = \nu(-A)$) and satisfies:
\begin{enumerate}[label=(L\arabic*)]
\item\label{L1} \textbf{Small-jump behavior (local stable type):}
  \[
    c_\gamma |z|^{-(1+\gamma)} \leq \frac{d\nu}{dz} \leq
    C_\gamma |z|^{-(1+\gamma)}, \qquad |z| \leq 1,
  \]
  for some $\gamma \in (1,2)$ and constants
  $0 < c_\gamma \leq C_\gamma$.
\item\label{L2} \textbf{Tail behavior:}
  $\int_{|z|>1} |z|^q\,\nu(dz) < \infty$ for some $q > 0$.
\end{enumerate}
The local stable index $\gamma$ governs the small-jump singularity
and is the primary parameter.  The tail condition~\ref{L2} is
deliberately left flexible: different results require different
values of $q$.  We distinguish two regimes:
\begin{enumerate}[label=(\alph*)]
\item \textbf{Tempered stable tails:}
  $\nu(dz) \sim |z|^{-1-\gamma}e^{-\lambda|z|}\,dz$ for
  $|z| > 1$, $\lambda > 0$.  Then~\ref{L2} holds for all
  $q < \infty$.  This is the regime in which our strongest
  results (sharpness, $L^2$ isometry) hold.
\item \textbf{Purely stable tails:}
  $\nu(dz) = c_\gamma|z|^{-1-\gamma}\,dz$ for all $z$.
  Then~\ref{L2} holds only for $q < \gamma$.  In this regime,
  several results require restricted moment conditions,
  which we flag explicitly.
\end{enumerate}
\end{assumption}

\begin{remark}[Moment notation]\label{rem:moments}
We write $\nu_r := \int_\R |z|^r\,\nu(dz)$ for the $r$-th
absolute moment of $\nu$, when finite.  By~\ref{L1}:
$\int_{|z|\leq 1}|z|^r\,\nu(dz) < \infty$ iff $r > \gamma$.
By~\ref{L2}: $\int_{|z|>1}|z|^r\,\nu(dz) < \infty$ iff
$r \leq q$ (or for all $r$ under tempered stable tails).
Thus $\nu_r < \infty$ iff $r > \gamma$ and $r \leq q$.

The \emph{core} defect decomposition
(\Cref{thm:defect-decomposition}) requires only:
$\int_{|z|\leq 1}|z|^{2p}\nu(dz) < \infty$ (automatic since
$2p > 2 > \gamma$) for the small-jump bound, and
$\int_{|z|>1}|z|^p\nu(dz) < \infty$ (automatic since
$p < \gamma$) for the large-jump bound.  In particular,
the core results hold under purely stable tails
without additional moment assumptions.

The \emph{full} moment $\nu_{2p} = \int_\R|z|^{2p}\nu(dz)$
is needed only for the sharpness proposition
(\Cref{thm:sharpness}), where exact norm identities over
all of $\R$ are claimed.  Under purely stable tails,
$\nu_{2p} = \infty$ (since $2p > 2 > \gamma$), and the
sharpness result does not apply.
\end{remark}

The case $\sigma > 0$ includes L\'evy processes with a Gaussian
component; $\sigma = 0$ is the pure-jump case.


\subsection{The Volterra--L\'evy process}

\begin{definition}[Volterra--L\'evy process]\label{def:VL}
Let $K : [0,T]^2 \to \R$ be a causal kernel ($K(t,s) = 0$ for
$s > t$). The \emph{Volterra--L\'evy process} is
\begin{equation}\label{eq:VL}
  V_t := \int_0^t K(t,s)\,dL_s
  = \sigma\!\int_0^t K(t,s)\,dW_s
  + \int_0^t\!\!\int_\R K(t,s)\,z\,\Ntilde(ds,dz).
\end{equation}
\end{definition}

\begin{assumption}[Kernel regularity]\label{ass:kernel}
The kernel $K$ satisfies:
\begin{enumerate}[label=(K\arabic*)]
\item\label{K1} \textbf{$L^\gamma$-integrability:}
  $\displaystyle\sup_{t \in [0,T]} \int_0^t |K(t,s)|^\gamma\,ds < \infty$.
  This ensures that $V_t$ is well-defined in $L^p$ for
  $p \in (1,\gamma)$.
\item\label{K2} \textbf{$L^2$-integrability (when $\sigma > 0$):}
  $\displaystyle\sup_{t \in [0,T]} \int_0^t K(t,s)^2\,ds < \infty$.
  This ensures the Gaussian part is well-defined in $L^2$.
\end{enumerate}
\end{assumption}

\begin{example}[Fractional kernel]\label{ex:fractional}
The \emph{fractional kernel} $K_H(t,s) = (t-s)^{H-1/2}$ for
$H \in (0,1)$ satisfies \ref{K1} iff $\gamma(H - 1/2) > -1$,
i.e., $H > 1/2 - 1/\gamma$.  For $\gamma > 1$, this holds for
all $H > 0$.  The Volterra--L\'evy process with $K_H$ and a
$\gamma$-stable driver is a \emph{rough stable process} with
roughness controlled by both $H$ and~$\gamma$.
\end{example}


\subsection{Semimartingale structure of the partial process}

The hedging analysis requires an It\^o formula. While $V_t$ as
a process in $t$ is generally \emph{not} a semimartingale (for
singular kernels), the \emph{partial process}
\begin{equation}\label{eq:partial}
  V_T(t) := \int_0^t K(T,s)\,dL_s, \qquad t \in [0,T],
\end{equation}
\emph{is} a semimartingale, since the kernel $K(T,\cdot)$ is
fixed (not time-varying) and the integral is taken with respect
to the semimartingale~$L$.

\begin{lemma}[Semimartingale decomposition of $V_T(\cdot)$]
\label{lem:SM-decomp}
The partial process $V_T(t)$ has the L\'evy--It\^o decomposition:
\begin{enumerate}[label=(\alph*)]
\item Continuous martingale part:
  $V_T^c(t) = \sigma\int_0^t K(T,s)\,dW_s$, with bracket
  $\langle V_T^c \rangle_t = \sigma^2\int_0^t K(T,s)^2\,ds$.
\item Jump part:
  $\Delta V_T(t) = K(T,t)\,\Delta L_t = K(T,t)\,\Delta J_t$,
  where $\Delta J_t$ are the jumps of
  $J_t = \int_0^t\!\int z\,\Ntilde(ds,dz)$.
\end{enumerate}
In particular, $V_T(T) = V_T$.
\end{lemma}

\begin{proof}
Since $K(T,\cdot)$ is deterministic, the stochastic integral
$\int_0^t K(T,s)\,dL_s$ is a semimartingale in $t$ by standard
theory~\cite{Applebaum2009}.  The decomposition follows from
linearity: $V_T(t) = V_T^c(t) + \int_0^t K(T,s)\,dJ_s$,
with $dJ_s = \int_\R z\,\Ntilde(ds,dz)$.
\end{proof}

\begin{remark}[Why $V_T(t)$ and not $V_t$]\label{rem:partial}
The process $V_t = \int_0^t K(t,s)\,dL_s$ is generally
\emph{not} a semimartingale when $K$ is singular (e.g.,
fractional kernel with $H < 1/2$).  The partial process
$V_T(t)$ freezes the terminal time $T$ and integrates
with respect to $L$, preserving the semimartingale property.
For terminal payoffs $f(V_T)$, the relevant object is
$f(V_T(T))$, and the It\^o formula can be applied to the
semimartingale $t \mapsto V_T(t)$.  This circumvents the
non-semimartingale difficulty entirely.
\end{remark}


%======================================================================
\section{The Combined Energy Space}\label{sec:energy}
%======================================================================

\subsection{Two-component energy space}

The L\'evy driver $L$ has both a continuous component (driven
by $W$) and a jump component (driven by $\Ntilde$).  Following
the Hilbert--Banach dichotomy established
in~\cite{Fontes2026a,Fontes2026d}, each component
has its natural energy space.

\begin{definition}[Combined energy space]\label{def:combined-energy}
The \emph{combined energy space} is the product Banach space
\begin{equation}\label{eq:HVL}
  \HH_{VL} := H_W \times \widetilde{H}_J,
\end{equation}
where:
\begin{enumerate}[label=(\roman*)]
\item The \emph{continuous energy space}:
  \begin{equation}\label{eq:HW}
    H_W := L^2\bigl([0,T] \times \Omega,\;
    \mathcal{P},\; \sigma^2\,ds \otimes d\PP\bigr),
  \end{equation}
  with norm
  $\|u\|_{H_W}^2 = \E\!\int_0^T u_t^2\,\sigma^2\,dt$
  and divergence
  $\delta_W(u) = \sigma\!\int_0^T u_t\,dW_t$.
  This is a Hilbert space.

\item The \emph{jump energy space}:
  \begin{equation}\label{eq:HJ}
    \widetilde{H}_J :=
    L^p\bigl([0,T] \times \R \times \Omega,\;
    \mathcal{P} \otimes \mathcal{B}(\R),\;
    dt \otimes \nu(dz) \otimes d\PP\bigr),
  \end{equation}
  with norm
  $\|h\|_{\widetilde{H}_J}^p = \E\!\int_0^T\!\!\int_\R
  |h(t,z)|^p\,\nu(dz)\,dt$
  and divergence
  $\widetilde{\delta}_J(h) = \int_0^T\!\!\int_\R
  h(t,z)\,\Ntilde(dt,dz)$.
  This is a Banach space (non-Hilbert for $p \neq 2$).
\end{enumerate}
The \emph{combined divergence} is
\begin{equation}\label{eq:delta-VL}
  \delta_{VL}(u,h) := \delta_W(u) + \widetilde{\delta}_J(h)
  = \sigma\!\int_0^T u_t\,dW_t
  + \int_0^T\!\!\int_\R h(t,z)\,\Ntilde(dt,dz).
\end{equation}
\end{definition}

\begin{remark}[Mixed exponent structure]\label{rem:mixed}
The continuous component lives in $L^2$ (Hilbert) while the jump
component lives in $L^p$ with $p < \gamma < 2$ (Banach).  We equip
$\HH_{VL}$ with the norm
$\|(u,h)\|_{\HH_{VL}} := \|u\|_{H_W} + \|h\|_{\widetilde{H}_J}$.
The combined divergence is bounded:
$\|\delta_{VL}(u,h)\|_{L^p} \leq
C_1\,\|u\|_{H_W} + C_2\,\|h\|_{\widetilde{H}_J}$
by the Brownian BDG inequality (for the $L^2$-bounded, hence
$L^p$-bounded, Brownian integral) and the L\'evy BDG inequality
(see~\cite{Dirksen2015}).

When $\sigma = 0$ (pure-jump case), $H_W$ is trivial and the
combined space reduces to $\widetilde{H}_J$ alone.
\end{remark}


\subsection{Combined operator derivative}

\begin{definition}[Combined operator derivative]\label{def:D-VL}
The \emph{combined operator derivative} is the Banach dual
\begin{equation}\label{eq:D-VL}
  D_{VL} := \delta_{VL}^* : L^q(\Omega) \to \HH_{VL}^*,
\end{equation}
defined by
$\ip{D_{VL} F}{(u,h)} = \E[F \cdot \delta_{VL}(u,h)]$
for $F \in L^q(\Omega)$, $(u,h) \in \HH_{VL}$.
\end{definition}

Since the combined divergence is the sum
$\delta_{VL}(u,h) = \delta_W(u) + \widetilde{\delta}_J(h)$,
linearity of the adjoint gives:

\begin{proposition}[Splitting of $D_{VL}$]\label{prop:D-split}
$D_{VL} F = (D_W F,\, \widetilde{D}_J F)$, where
\begin{align}
  \ip{D_W F}{u}_{H_W^*, H_W}
  &= \E\!\left[F \cdot \sigma\!\int_0^T u_t\,dW_t\right],
  \label{eq:DW}\\
  \ip{\widetilde{D}_J F}{h}_{\widetilde{H}_J^*, \widetilde{H}_J}
  &= \E\!\left[F \cdot \int_0^T\!\!\int_\R
  h(t,z)\,\Ntilde(dt,dz)\right].
  \label{eq:DJ}
\end{align}
The component $D_W$ is the Hilbert operator derivative
of~\cite{Fontes2026a}, and $\widetilde{D}_J$ is the
two-parameter Banach derivative (\Cref{def:D-VL}).
\end{proposition}


%======================================================================

\subsection{Why not Malliavin calculus?}\label{sec:why-not-malliavin}

The combined operator derivative $D_{VL}$ is defined by a single
mechanism---duality with the stochastic integral---applied to the
product space $\mathcal{H}_{VL} = H_W \times \widetilde{H}_J$.
The splitting $D_{VL} F = (D_W F, \widetilde{D}_J F)$ is
\emph{inherited} from the direct-sum structure of the energy
space, not constructed from separate calculi.

This stands in contrast to the Malliavin-theoretic approach,
which requires two distinct tools: the Malliavin derivative on
Wiener space (for the Brownian component) and the Picard or
Nualart--Vives difference operator on Poisson space (for the
jumps).  These are defined by different mechanisms (Gaussian
integration by parts vs.\ finite differences in the jump
configuration) and must be stitched together manually when both
are present.  The kernel convolution interacting with both
components must be managed through separate computations.

The operator derivative circumvents this entirely.  Because
$D_{VL} = \delta_{VL}^*$ is defined by the adjoint identity
$\langle D_{VL} F, (u,h) \rangle = \mathbb{E}[F \cdot
\delta_{VL}(u,h)]$, the kernel-weighted defect
$\Gamma_V^f$---the Taylor remainder of the composition
$f(V_T(t\!-\!) + K(T,t)z) - f(V_T(t\!-\!))
- f'(V_T(t\!-\!))K(T,t)z$---emerges as a single algebraic
object in $\widetilde{H}_J$.  In the language
of~\cite{Fontes2026d}, the defect is the non-Gaussian content:
$D$ satisfies the Leibniz rule if and only if $\nu = 0$
(no jumps), i.e., if and only if the driver is Gaussian.
The failure of Leibniz, measured by $\Gamma_V^f$, is therefore
exactly the component of risk that Gaussian tools cannot see.


\section{The Volterra--L\'evy It\^o Formula}\label{sec:ito}
%======================================================================

This section establishes the decomposition of $f(V_T)$ into
hedgeable and unhedgeable components.  The unhedgeable
part---the kernel-weighted Leibniz defect---is the exact
content that Malliavin calculus cannot access: it measures
the nonlinear response of the payoff to jumps propagated
through the kernel, which is the Leibniz failure of the
combined operator derivative.  The construction applies the
classical It\^o formula for semimartingales to the partial
process $V_T(t)$, then decomposes the result using the
operator framework.

\begin{theorem}[Volterra--L\'evy It\^o formula]\label{thm:VL-ito}
Let $f \in C_b^3(\R)$ and $V_T = \int_0^T K(T,s)\,dL_s$ under
Assumptions~\ref{ass:nu} and~\ref{ass:kernel}.  Then
\begin{equation}\label{eq:VL-ito}
  f(V_T) - f(0)
  = \underbrace{\int_0^T \alpha_t^f\,\sigma\,dW_t}_{I_W^f}
  + \underbrace{\int_0^T\!\!\int_\R h_f^K(t,z)\,\Ntilde(dt,dz)}_{I_J^f}
  + \underbrace{A_T^f}_{\text{drift}},
\end{equation}
where, writing $V_T(t\!-\!) := \int_0^{t-} K(T,s)\,dL_s$:
\begin{align}
  \alpha_t^f &:= f'(V_T(t\!-\!))\,K(T,t),
  \label{eq:alpha}\\[4pt]
  h_f^K(t,z) &:= f\bigl(V_T(t\!-\!) + K(T,t)\,z\bigr)
  - f(V_T(t\!-\!)),
  \label{eq:hfK}\\[4pt]
  g_f^K(t,z) &:= h_f^K(t,z) - f'(V_T(t\!-\!))\,K(T,t)\,z
  \label{eq:gfK-thm}
\end{align}
(the Taylor remainder), and the drift is
\begin{equation}\label{eq:drift}
  A_T^f := \tfrac{1}{2}\sigma^2\!\int_0^T
  f''(V_T(t\!-\!))\,K(T,t)^2\,dt
  + \int_0^T\!\!\int_\R g_f^K(t,z)\,\nu(dz)\,dt.
\end{equation}
The second integral converges absolutely.
For $|z| \leq 1$:
$|g_f^K(t,z)| \leq \frac{1}{2}\|f''\|_\infty K(T,t)^2 z^2$
by Taylor's theorem, and $\int_{|z|\leq 1}|z|^2\,\nu(dz) < \infty$
since $\gamma < 2$.
For $|z| > 1$: $|g_f^K(t,z)| \leq 2\|f\|_\infty
+ \|f'\|_\infty|K(T,t)||z|$, and
$\int_{|z|>1}|z|\,\nu(dz) < \infty$ since $\gamma > 1$
(the local stable assumption~\ref{L1} gives
$\int_1^\infty z \cdot z^{-(1+\gamma)}\,dz < \infty$
for $\gamma > 1$, and the tail condition~\ref{L2} controls
the large-$|z|$ remainder).
\end{theorem}

\begin{proof}
Apply the It\^o formula for semimartingales with jumps
(\cite{Applebaum2009}, Theorem~4.4.7) to the
semimartingale $V_T(t)$
(\Cref{lem:SM-decomp}):
\begin{align*}
  f(V_T(T)) &= f(V_T(0))
  + \int_0^T f'(V_T(t\!-\!))\,dV_T(t)
  + \tfrac{1}{2}\int_0^T f''(V_T(t\!-\!))\,d\langle V_T^c\rangle_t\\
  &\quad+ \sum_{0 < s \leq T}\left[f(V_T(s)) - f(V_T(s\!-\!))
  - f'(V_T(s\!-\!))\,\Delta V_T(s)\right].
\end{align*}
Since $V_T(0) = 0$, $V_T(T) = V_T$,
$dV_T(t) = K(T,t)\,dL_t$, and
$\Delta V_T(s) = K(T,s)\,\Delta L_s$:

\textbf{Term 1.} The first-order stochastic integral.
Since $dL_t = \sigma\,dW_t + \int_\R z\,\Ntilde(dt,dz)$
(the driver is defined with the compensated measure, so
there is no drift):
\begin{align*}
  \int_0^T f'(V_T(t\!-\!))\,K(T,t)\,dL_t
  &= \int_0^T f'(V_T(t\!-\!))\,K(T,t)\,\sigma\,dW_t\\
  &\quad+ \int_0^T\!\!\int_\R f'(V_T(t\!-\!))\,K(T,t)\,z\,
  \Ntilde(dt,dz).
\end{align*}

\textbf{Term 2.} The quadratic variation:
$\frac{1}{2}\int_0^T f''(V_T(t\!-\!))\,K(T,t)^2\,\sigma^2\,dt$.

\textbf{Term 3.} The jump sum.  Writing
$g_f^K$ for the Taylor remainder~\eqref{eq:gfK-thm}:
\begin{align*}
  &\sum_{s \leq T}\bigl[f(V_T(s\!-\!) + K(T,s)\Delta L_s)
  - f(V_T(s\!-\!)) - f'(V_T(s\!-\!))K(T,s)\Delta L_s\bigr]\\
  &= \int_0^T\!\!\int_\R g_f^K(t,z)\,N(dt,dz)\\
  &= \int_0^T\!\!\int_\R g_f^K(t,z)\,\Ntilde(dt,dz)
  + \int_0^T\!\!\int_\R g_f^K(t,z)\,\nu(dz)\,dt.
\end{align*}

\textbf{Collecting terms.}  The Brownian martingale part is
$I_W^f = \int_0^T f'(V_T(t\!-\!))\,K(T,t)\,\sigma\,dW_t$
with $\alpha_t^f = f'(V_T(t\!-\!))\,K(T,t)$.

The compensated Poisson integral combines the first-order term
and the defect:
\begin{align*}
  I_J^f &= \int_0^T\!\!\int_\R
  \bigl[f'(V_T(t\!-\!))K(T,t)z + g_f^K(t,z)\bigr]\,
  \Ntilde(dt,dz)\\
  &= \int_0^T\!\!\int_\R h_f^K(t,z)\,\Ntilde(dt,dz),
\end{align*}
where $h_f^K(t,z) = f(V_T(t\!-\!) + K(T,t)z) - f(V_T(t\!-\!))$
is the full finite difference~\eqref{eq:hfK}.

The drift collects the quadratic variation and the compensated
jump terms, yielding~\eqref{eq:drift}.
\end{proof}


%----------------------------------------------------------------------
\subsection{The kernel-weighted defect decomposition}
%----------------------------------------------------------------------

The compensated Poisson integrand~\eqref{eq:hfK} decomposes into
a \emph{linear part} (first-order response to jumps) and a
\emph{nonlinear part} (the Taylor remainder):
\begin{equation}\label{eq:hfK-decomp}
  h_f^K(t,z)
  = \underbrace{f'(V_T(t\!-\!))\,K(T,t)\,z}_{\text{linear in }z}
  + \underbrace{g_f^K(t,z)}_{\text{Leibniz defect integrand}}.
\end{equation}

The integrability of these two components in $\widetilde{H}_J$
differs qualitatively.

\begin{theorem}[Kernel-weighted defect decomposition]
\label{thm:defect-decomposition}
Under Assumptions~\ref{ass:nu} and~\ref{ass:kernel}, let
$f \in C_b^3(\R)$ and $p \in (1,\gamma)$.  Assume further that
$\int_0^T |K(T,t)|^{2p}\,dt < \infty$ (automatic for bounded
kernels; for the fractional kernel, this is \Cref{ass:Hp}).
Then:
\begin{enumerate}[label=(\alph*)]
\item \textbf{$\nu$-regularity of the defect.}
  The kernel-weighted Taylor remainder $g_f^K$ lies in
  $\widetilde{H}_J$:
  \begin{equation}\label{eq:gfK-bound}
    \|g_f^K\|_{\widetilde{H}_J}^p
    \leq C_{\mathrm{small}}\!\int_0^T\!\!|K(T,t)|^{2p}\,dt
    \;\cdot\;\int_{|z|\leq 1}|z|^{2p}\,\nu(dz)
    + C_{\mathrm{large}}\,T\!\int_{|z|>1}\!(1+|z|^p)\,\nu(dz),
  \end{equation}
  where $C_{\mathrm{small}} = C(p)\,\|f''\|_\infty^p/2^p$
  and $C_{\mathrm{large}}$ depends on $\|f\|_\infty$,
  $\|f'\|_\infty$.

\item \textbf{$\nu$-singularity of the linear part.}
  The integrand $f'(V_T(t\!-\!))\,K(T,t)\,z$ does not belong to
  $\widetilde{H}_J$: its norm involves
  $\int_\R |z|^p\,\nu(dz) = \infty$ (since $p < \gamma < 2$ and
  $\nu$ is $\gamma$-stable near zero).

\item \textbf{Two-regime decomposition.}
  The compensated Poisson integral $I_J^f$ decomposes as:
  \begin{equation}\label{eq:two-regime-VL}
    I_J^f
    = \underbrace{\int_0^T\!\!\int_\R
      f'(V_T(t\!-\!))K(T,t)z\,\Ntilde(dt,dz)}_{\text{%
      $\nu$-singular (linear in $z$)}}
    + \underbrace{\int_0^T\!\!\int_\R
      g_f^K(t,z)\,\Ntilde(dt,dz)}_{\text{%
      $\nu$-regular (defect)}}.
  \end{equation}
  The $\nu$-regular part is the kernel-weighted Leibniz defect:
  \begin{equation}\label{eq:GammaVf}
    \Gamma_V^f := \widetilde{\delta}_J(g_f^K)
    = \int_0^T\!\!\int_\R g_f^K(t,z)\,\Ntilde(dt,dz).
  \end{equation}
\end{enumerate}
\end{theorem}

\begin{proof}
(a) For $|z| \leq 1$: by Taylor's theorem,
\[
  |g_f^K(t,z)| = |f(V_T(t\!-\!) + K(T,t)z)
  - f(V_T(t\!-\!)) - f'(V_T(t\!-\!))K(T,t)z|
  \leq \tfrac{1}{2}\|f''\|_\infty |K(T,t)|^2 |z|^2.
\]
Raising to the $p$-th power and integrating:
\begin{align*}
  \E\!\int_0^T\!\!\int_{|z|\leq 1}|g_f^K|^p\,\nu(dz)\,dt
  &\leq \frac{\|f''\|_\infty^p}{2^p}
  \int_0^T |K(T,t)|^{2p}\,dt
  \cdot\int_{|z|\leq 1}|z|^{2p}\,\nu(dz).
\end{align*}
The time integral converges by~\ref{K1} (or verified directly for
fractional kernels; see \Cref{lem:kernel-integral}).  The jump
integral converges since $2p > \gamma$:
$\int_{|z|\leq 1}|z|^{2p}\,\nu(dz)
\leq C_\gamma\int_0^1 z^{2p - 1 - \gamma}\,dz
= C_\gamma/(2p - \gamma) < \infty$.

For $|z| > 1$: $|g_f^K(t,z)| \leq 2\|f\|_\infty
+ \|f'\|_\infty|K(T,t)||z|$, so $|g_f^K|^p \leq
C(1 + |K(T,t)|^p|z|^p)$.  Since $p < \gamma$,
$\int_{|z|>1}|z|^p\,\nu(dz) < \infty$, and
$\int_0^T |K(T,t)|^p\,dt < \infty$ by~\ref{K1}.

(b) $\int_{|z|\leq 1}|z|^p\,\nu_\gamma(dz)
= 2c_\gamma\int_0^1 z^{p-1-\gamma}\,dz = \infty$
since $p - \gamma < 0$.  By the uniqueness of the compensated
Poisson representation, $f'(V_T(t\!-\!))K(T,t)z$ is the only
possible two-parameter integrand for $\delta_L(f'(V_T(\cdot\!-\!))
K(T,\cdot))$, and it does not belong to $\widetilde{H}_J$.

(c) Combine (a), (b), and the It\^o formula~\eqref{eq:VL-ito}.
\end{proof}

\begin{remark}[Interpretation: what is the defect?]
\label{rem:defect-interpret}
The kernel-weighted Leibniz defect $\Gamma_V^f$ encodes the
\emph{nonlinear response of $f(V_T)$ to jumps in the driver $L$,
propagated through the kernel $K$}.  A jump of size $z$ at time $t$
shifts $V_T$ by $K(T,t)\,z$; the defect integrand $g_f^K(t,z)$
measures the second-order effect of this shift.  For smooth $f$,
$g_f^K(t,z) \approx \frac{1}{2}f''(V_T(t\!-\!))K(T,t)^2 z^2$
for small $z$---it is the ``gamma times jump-squared'' contribution,
weighted by the kernel.

This is the stochastic-calculus analog of the Leibniz defect
in~\cite{Fontes2026d}: the failure of $D_L$ to satisfy the
product rule generates a correction term that is intrinsically
tied to the jump structure.  The Volterra twist is that the kernel
$K(T,t)$ amplifies or attenuates each jump's contribution depending
on when it occurs.
\end{remark}


\subsection{Extension to polynomial-growth payoffs}
\label{sec:poly-growth}

The It\^o formula (\Cref{thm:VL-ito}) and the defect
decomposition (\Cref{thm:defect-decomposition}) are stated
for $f \in C_b^3(\R)$.  Many payoffs of interest---including
$f(x) = x^2$ (variance swaps) and $f(x) = e^x$
(exponential volatility)---do not satisfy this assumption.
We now extend the main results to functions with polynomial
growth.

\begin{definition}[Polynomial growth class]\label{def:poly-growth}
For $m \geq 0$, define $C_{\mathrm{pol}}^3(\R,m)$ as the
class of $f \in C^3(\R)$ satisfying
\[
  |f^{(k)}(x)| \leq C_f\,(1 + |x|^m)
  \qquad \text{for }k = 0,1,2,3,
\]
for some constant $C_f > 0$.  The case $m = 0$ recovers
$C_b^3(\R)$.
\end{definition}

\begin{assumption}[Moment condition for polynomial payoffs]
\label{ass:moments-poly}
The Volterra--L\'evy process $V_T$ satisfies
$\E[|V_T|^r] < \infty$ for all $r$ up to a specified order,
and the tail condition~\ref{L2} holds with $q$ large
enough.  Specifically, for $f \in C_{\mathrm{pol}}^3(\R,m)$
and working exponent $p$, we require:
\begin{enumerate}[label=(M\arabic*)]
\item\label{M1} $\E[|V_T|^{(m+2)p}] < \infty$
  (moment condition on the Volterra process);
\item\label{M2} $\int_{|z|>1}|z|^{(m+1)p}\,\nu(dz) < \infty$
  (tail condition: \ref{L2} with $q = (m+1)p$).
\end{enumerate}
\end{assumption}

\begin{remark}[When do these moments exist?]\label{rem:VT-moments}
For a $\gamma$-stable L\'evy driver $L$ and bounded kernel $K$,
$V_T = \int_0^T K(T,s)\,dL_s$ inherits the moment structure of
$L$: under tempered stable tails, $\E[|V_T|^r] < \infty$ for
all $r < \infty$; under purely stable tails, $\E[|V_T|^r] < \infty$
only for $r < \gamma$.  Thus~\ref{M1} is automatic under tempered
stable tails and severely restrictive under purely stable tails
(requiring $(m+2)p < \gamma$, which conflicts with $p > 1$ for
$m \geq 1$ and $\gamma < 2$).  The polynomial growth extension
is therefore primarily applicable under tempered stable or
finite-activity jump assumptions.
\end{remark}

\begin{corollary}[Volterra--L\'evy It\^o formula for polynomial
growth]\label{cor:VL-ito-poly}
Let $f \in C_{\mathrm{pol}}^3(\R,m)$ and assume
\Cref{ass:moments-poly}.  Then the decomposition~\eqref{eq:VL-ito}
holds, with each term well-defined in $L^p$:
\begin{enumerate}[label=(\alph*)]
\item The Brownian integral: $\E[|\int_0^T
  f'(V_T(t\!-\!))K(T,t)\sigma\,dW_t|^p] < \infty$, since
  $|f'(V_T(t\!-\!))| \leq C(1+|V_T(t\!-\!)|^m)$ and the
  BDG inequality gives $\E[|I_W^f|^p] \leq
  C\,\E[\int_0^T(1+|V_T(t)|^m)^2 K(T,t)^2\,dt]^{p/2}$,
  which is finite by~\ref{M1}.
\item The compensated Poisson integral: the integrand
  $h_f^K(t,z) = f(V_T(t\!-\!) + K(T,t)z) - f(V_T(t\!-\!))$
  satisfies $|h_f^K(t,z)| \leq
  C(1 + |V_T(t\!-\!)|^m + |K(T,t)z|^{m+1})$ by the
  mean value theorem.  The $L^p(dt \otimes \nu \otimes dP)$
  integrability follows from~\ref{M1} and~\ref{M2}.
\item The drift: $|A_T^f| \leq C\int_0^T(1+|V_T(t)|^m)
  K(T,t)^2\,dt + C\int_0^T\!\int|h_f^K|\,\nu(dz)\,dt$,
  which is finite under the same conditions.
\end{enumerate}
\end{corollary}

\begin{proof}
The It\^o formula for semimartingales with jumps
(\cite{Applebaum2009}, Theorem~4.4.7) requires only that
$f \in C^2$ and that the relevant stochastic integrals
and jump sums converge.  The $C_b^3$ assumption in
\Cref{thm:VL-ito} was used for the uniform bound
$|g_f^K| \leq \frac{1}{2}\|f''\|_\infty K^2 z^2$ in the
defect decomposition; for polynomial-growth $f$, this
bound becomes $|g_f^K(t,z)| \leq
C(1+|V_T(t\!-\!)|^m)\,K(T,t)^2\,z^2$ for $|z| \leq 1$
(by Taylor's theorem applied to $f$ at $V_T(t\!-\!)$),
and $|g_f^K(t,z)| \leq C(1+|V_T(t\!-\!)|^{m+1}
+ |K(T,t)z|^{m+1})$ for $|z| > 1$.

The $\widetilde{H}_J$-norm of $g_f^K$ is then bounded by:
\begin{align*}
  \|g_f^K\|_{\widetilde{H}_J}^p
  &\leq C^p\,\E\!\int_0^T\!(1+|V_T(t)|^m)^p K(T,t)^{2p}\,dt
  \;\cdot\;\int_{|z|\leq 1}|z|^{2p}\,\nu(dz)\\
  &\quad+ C^p\,\E\!\int_0^T\!\!\int_{|z|>1}\!
  \bigl(1+|V_T(t)|^{(m+1)p}
  + |K(T,t)|^{(m+1)p}\,|z|^{(m+1)p}\bigr)\,\nu(dz)\,dt.
\end{align*}
The first line is finite by~\ref{M1} (which gives
$\E[(1+|V_T(t)|^m)^p] < \infty$) and the kernel integrability
condition $\int_0^T K(T,t)^{2p}\,dt < \infty$ (which must be
assumed; for fractional kernels, this is $H > (p-1)/(2p)$,
see \Cref{ass:Hp} in Section~\ref{sec:scaling}).
The second line is finite by~\ref{M1} and~\ref{M2}.
\end{proof}

\begin{remark}[The quadratic payoff revisited]
\label{rem:quadratic-revisited}
For $f(x) = x^2$: $m = 2$, $f''(x) = 2$, $f'''(x) = 0$,
so $f \in C_{\mathrm{pol}}^3(\R, 2)$ with the polynomial
growth arising only in $f$ and $f'$, not in the derivatives
used for the Taylor bound.  In fact, the defect integrand
$g_f^K(t,z) = K(T,t)^2 z^2$ is \emph{independent of
$V_T(t\!-\!)$}, so the moment condition~\ref{M1} is not
needed for the defect norm computation---only~\ref{M2}
(requiring $\nu_{2p} < \infty$) is needed.  This is why
\Cref{thm:sharpness} requires only tail conditions on $\nu$,
not moment conditions on $V_T$.
\end{remark}


%======================================================================
\section{The Fundamental Pricing Constraint}\label{sec:pricing}
%======================================================================

The results of the preceding sections---the Volterra--L\'evy
It\^o formula, the defect decomposition, and the
$\nu$-singular/$\nu$-regular splitting---describe the
\emph{structure} of functionals of $V_T$.  We now turn to
the question they are built to answer: what does risk-neutral
pricing mean in a market with memory and jumps?

The answer is the two-parameter Volterra risk-neutral
constraint (VRNC): in a market with memory and jumps, the
risk-neutral condition is a Volterra integral equation.  The
central finding is that the kernel~$K$ enters the pricing
constraint itself: the Brownian and jump risk premia are
coupled through a Volterra integral equation, and cannot
be chosen independently.  We also establish that the
roughness parameter~$H$ is a structural invariant of the
market, identical under the physical and pricing measures.

\subsection{Asset dynamics with memory and jumps}

Consider an asset driven by both the L\'evy driver $L$ directly
and by a memory factor:
\begin{equation}\label{eq:asset-VL}
  \frac{dS_t}{S_t} = \mu_t\,dt + \lambda_t\,dL_t
  + \beta_t\,dZ_t,
\end{equation}
where $Z_t = \int_0^t K(t,s)\,dL_s$ is the Volterra observable,
$\lambda_t$ is the direct (``fresh'') exposure, and $\beta_t$ is
the memory-factor loading.  Both $\lambda$ and $\beta$ are bounded
predictable processes.

\begin{remark}[Semimartingale scope]\label{rem:VRNC-scope}
The dynamics~\eqref{eq:asset-VL} require $Z$ to have
well-defined stochastic differentials, i.e., $Z$ must be a
semimartingale.  For regular kernels (e.g., exponential
$K(t,s) = e^{-\kappa(t-s)}$), $Z$ is a semimartingale and
$dZ_t$ is well-defined.  For singular kernels (fractional
$K_H$ with $H < 1/2$), $Z$ is generally not a semimartingale,
and the ``differential form''~\eqref{eq:asset-VL} is
interpreted at the integrated level:
$\log S_t = \log S_0 + \int_0^t(\mu_s - \frac{1}{2}\sigma_s^2)\,ds
+ \lambda \cdot L_t + \beta \cdot Z_t + \ldots$, where
$\lambda \cdot L$ and $\beta \cdot Z$ are Volterra integrals.
The integral form of the VRNC~\eqref{eq:VRNC-int} is then the
primary constraint.
\end{remark}


\subsection{L\'evy--Girsanov measure change}

An equivalent measure $\QQ$ is defined by the Girsanov change on
the two noise components of $L = \sigma_L W + J$:
\begin{enumerate}[label=(\roman*)]
\item \textbf{Continuous:} $W_t^\QQ := W_t + \int_0^t
  \theta_s^W\,ds$, with $\theta^W$ a bounded predictable process
  satisfying the Novikov condition.
\item \textbf{Jump:} The compensated measure becomes
  $\Ntilde^\QQ(dt,dz) := N(dt,dz) - e^{\theta_t^J(z)}\nu(dz)\,dt$,
  with $\theta^J : [0,T] \times \R \to \R$ predictable in $t$,
  measurable in $z$, satisfying
  $\int_\R(e^{\theta^J(z)} - 1 - \theta^J(z))\nu(dz) < \infty$.
\end{enumerate}
The existence and form of the $\QQ$-density follow from
the Girsanov theorem (for the Brownian component) and
the L\'evy--Girsanov theorem (for the jump component);
see~\cite{ContTankov2004}, Theorem~9.12, for the precise
integrability conditions and the Radon--Nikodym density.
We assume throughout that $(\theta^W, \theta^J)$ satisfy
these conditions.


\subsection{Drift functional for Volterra--L\'evy observables}

\begin{theorem}[Two-parameter drift functional]
\label{thm:drift-VL}
Under the L\'evy--Girsanov measure change $(\theta^W, \theta^J)$,
the Volterra observable $Z_t = \int_0^t K(t,s)\,dL_s$ decomposes
under $\QQ$ as:
\begin{equation}\label{eq:Z-decomp-VL}
  Z_t = Z_t^\QQ + B_t^\theta,
\end{equation}
where $Z_t^\QQ = \int_0^t K(t,s)\,dL_s^\QQ$ is the
$\QQ$-observable and the drift functional is
\begin{equation}\label{eq:drift-functional-VL}
  B_t^\theta
  = \int_0^t K(t,s)\,\Psi_s\,ds,
\end{equation}
with the \emph{effective drift rate}
\begin{equation}\label{eq:Psi-def}
  \Psi_s := -\sigma_L\,\theta_s^W
  + \int_\R z\bigl(e^{\theta_s^J(z)} - 1\bigr)\,\nu(dz).
\end{equation}
The first term is the continuous Girsanov drift; the second is
the jump-induced drift.
\end{theorem}

\begin{proof}
Under $\QQ$, the L\'evy driver decomposes as:
\begin{align*}
  dL_t &= \sigma_L\,dW_t + \int_\R z\,\Ntilde(dt,dz)\\
  &= \sigma_L\,dW_t^\QQ - \sigma_L\theta_t^W\,dt
  + \int_\R z\,\Ntilde^\QQ(dt,dz)
  + \int_\R z\bigl(e^{\theta_t^J(z)} - 1\bigr)\nu(dz)\,dt\\
  &= dL_t^\QQ + \Psi_t\,dt.
\end{align*}
Substituting into the Volterra integral:
\begin{align*}
  Z_t &= \int_0^t K(t,s)\,dL_s
  = \int_0^t K(t,s)\,dL_s^\QQ + \int_0^t K(t,s)\,\Psi_s\,ds
  = Z_t^\QQ + B_t^\theta.
\end{align*}
The drift functional $B_t^\theta = \int_0^t K(t,s)\Psi_s\,ds$ is a
Volterra integral of $\Psi$ against the kernel---exactly the same
structural form as the continuous case
(\Cref{cor:VRNC-recovery}), with $\Psi$ replacing the
continuous-only drift rate $-A_s\theta_s$.
\end{proof}

\begin{remark}[Structure of $\Psi$]\label{rem:Psi-structure}
The effective drift rate $\Psi_s$ aggregates the continuous and jump
market prices of risk into a single scalar at each time~$s$.  The
jump contribution $\int z(e^{\theta^J(z)} - 1)\nu(dz)$ is the
mean of the L\'evy measure under the tilted weight
$e^{\theta^J(z)}$, minus the original mean (which is zero by
symmetry of $\nu$).  This integral is well-defined when
$\theta^J$ satisfies the integrability condition in the
L\'evy--Girsanov theorem.

For the Esscher transform $\theta^J(z) = \vartheta z$ (a common
choice in financial applications),
$\Psi_s = -\sigma_L\theta_s^W + \int z(e^{\vartheta_s z} - 1)\nu(dz)$.
For symmetric $\nu$, the integral simplifies to
$\int z(e^{\vartheta z} - 1)\nu(dz)
= 2\int_0^\infty z\sinh(\vartheta z)\nu(dz)$,
which is monotone in $\vartheta$.
\end{remark}


\subsection{The two-parameter VRNC}

\begin{theorem}[Two-parameter Volterra risk-neutral constraint]
\label{thm:VRNC-VL}
Under the asset dynamics~\eqref{eq:asset-VL} with memory factor
$Z_t = \int_0^t K(t,s)\,dL_s$, suppose $Z$ is a semimartingale
(this holds for regular kernels, e.g.\ exponential
$K(t,s) = e^{-\kappa(t-s)}$; for singular kernels see
\Cref{rem:VRNC-singular}).  An equivalent measure $\QQ$
defined by $(\theta^W, \theta^J)$ satisfying the
Girsanov and L\'evy--Girsanov integrability conditions
(see~\cite{ContTankov2004}, Theorem~9.12)
makes the discounted stock
$\tilde{S}$ a $\QQ$-local martingale only if the
effective drift rate $\Psi$ satisfies:
\begin{enumerate}[label=(\alph*)]
\item \textbf{Differential form} (for regular kernels with
  $\partial_t K$ existing):
  \begin{equation}\label{eq:VRNC-diff}
    \mu_t - r
    = -\bigl[\lambda_t + \beta_t\,K(t,t)\bigr]\,\Psi_t
    - \beta_t\!\int_0^t \partial_t K(t,s)\,\Psi_s\,ds.
  \end{equation}

\item \textbf{Integral form} (for all kernels, including singular):
  \begin{equation}\label{eq:VRNC-int}
    \int_0^t (\mu_s - r)\,ds
    = -\int_0^t \lambda_s\,\Psi_s\,ds
    - \int_0^t \beta_s\,\dot{B}_s^\theta\,ds,
  \end{equation}
  where $\dot{B}_s^\theta = \frac{d}{ds}\int_0^s K(s,u)\Psi_u\,du$
  when the derivative exists.  For singular kernels,
  the integral form is stated directly as the integrated
  version of~\eqref{eq:RN-condition}.
\end{enumerate}
Both forms are Volterra integral equations for $\Psi$
(of the second kind in the differential form, of the first kind
in the integral form).  They are structurally identical to the
continuous case (recovered in \Cref{cor:VRNC-recovery}), with the effective drift
rate $\Psi_t = -\sigma_L\theta_t^W
+ \int z(e^{\theta_t^J(z)} - 1)\nu(dz)$ replacing the
continuous-only $-A_t\theta_t$.
\end{theorem}

\begin{proof}
Under $\QQ$, the asset dynamics~\eqref{eq:asset-VL} become:
\begin{align*}
  \frac{dS_t}{S_t}
  &= \mu_t\,dt + \lambda_t\,(dL_t^\QQ + \Psi_t\,dt)
  + \beta_t\,(dZ_t^\QQ + dB_t^\theta)\\
  &= [\mu_t + \lambda_t\Psi_t + \beta_t\dot{B}_t^\theta]\,dt
  + (\text{$\QQ$-martingale terms}).
\end{align*}
For $\tilde{S}_t = e^{-rt}S_t$ to be a $\QQ$-local martingale,
the $dt$-coefficient must vanish:
\begin{equation}\label{eq:RN-condition}
  \mu_t - r + \lambda_t\,\Psi_t + \beta_t\,\dot{B}_t^\theta = 0.
\end{equation}

\textbf{Differential form.}  For regular kernels, the Leibniz rule
gives:
$\dot{B}_t^\theta = K(t,t)\,\Psi_t + \int_0^t \partial_t K(t,s)\,
\Psi_s\,ds$.
Substituting into~\eqref{eq:RN-condition}:
$\mu_t - r = -\lambda_t\Psi_t - \beta_t[K(t,t)\Psi_t
+ \int_0^t\partial_t K\,\Psi_s\,ds]$,
yielding~\eqref{eq:VRNC-diff}.

\textbf{Integral form.}  Integrate~\eqref{eq:RN-condition} over
$[0,t]$ directly:
$\int_0^t(\mu_s - r)\,ds = -\int_0^t \lambda_s\Psi_s\,ds
- \int_0^t \beta_s\dot{B}_s^\theta\,ds$.
For singular kernels where $\dot{B}^\theta$ does not exist
classically, this integrated form is the defining relation.
\end{proof}

\begin{corollary}[Recovery of the continuous VRNC]
\label{cor:VRNC-recovery}
When $\nu = 0$ (no jumps), $\theta^J$ is absent,
$\Psi_t = -\sigma_L\theta_t^W$, and~\eqref{eq:VRNC-diff}
reduces to the continuous
VRNC:
\[
  \mu_t - r = [\lambda_t + \beta_t K(t,t)]\sigma_L\theta_t^W
  + \beta_t\int_0^t \partial_t K(t,s)\,\sigma_L\theta_s^W\,ds.
\]
\end{corollary}

\begin{corollary}[Markov recovery]\label{cor:VRNC-markov}
When $K(t,s) = \mathbf{1}_{\{s \leq t\}}$ (Markov kernel),
$K(t,t) = 1$ and $\partial_t K = 0$.  The VRNC reduces to
the \emph{pointwise} constraint:
\[
  \mu_t - r = -(\lambda_t + \beta_t)\,\Psi_t,
\]
which determines $\Psi_t$ algebraically at each $t$ (no integral
equation).  In the continuous case ($\Psi_t = -\sigma_L\theta_t^W$),
this becomes $\mu_t - r = (\lambda_t + \beta_t)\sigma_L\theta_t^W$,
the classical drift cancellation condition.
\end{corollary}


\subsection{The measure-selection degrees of freedom}

\begin{proposition}[Decomposition of the pricing freedom]
\label{prop:freedom}
The VRNC (\Cref{thm:VRNC-VL}) determines the effective drift
rate $\Psi$ uniquely (when the Volterra equation has a unique
solution).  However, $\Psi$ does not determine $(\theta^W,
\theta^J)$ uniquely: the decomposition
\begin{equation}\label{eq:Psi-decomp}
  \Psi_t = -\sigma_L\,\theta_t^W + \Phi_t^J,
  \qquad
  \Phi_t^J := \int_\R z\bigl(e^{\theta_t^J(z)} - 1\bigr)\,\nu(dz),
\end{equation}
has one equation in the unknowns $(\theta_t^W, \theta_t^J(\cdot))$.
The remaining degrees of freedom are:
\begin{enumerate}[label=(\alph*)]
\item \textbf{The continuous--jump split:} For any choice of
  $\Phi^J$ (equivalently, $\theta^J$), the continuous component is
  determined: $\theta_t^W = (\Phi_t^J - \Psi_t)/\sigma_L$.
  Different choices of $\Phi^J$ correspond to different
  allocations of the effective drift between the Brownian and
  jump channels.

\item \textbf{The jump-measure shape:} Even given $\Phi^J$
  (the first $z$-moment of the tilted measure), the function
  $\theta^J(z)$ is not determined: many L\'evy measure
  changes $e^{\theta^J(z)}\nu(dz)$ can produce the same
  $\Phi^J = \int ze^{\theta^J(z)}\nu(dz)$.  This is the
  infinite-dimensional freedom of the jump-risk pricing
  kernel.
\end{enumerate}
The total pricing freedom is therefore: one functional
degree of freedom (the choice of $\Phi^J_t$ at each $t$)
plus an infinite-dimensional degree per time point (the
shape of $\theta^J_t(\cdot)$ conditional on its
$z$-weighted integral being $\Phi_t^J$).
\end{proposition}

\begin{proof}
The VRNC is a Volterra equation for $\Psi$, which (under
standard conditions on $\lambda, \beta, K$---see,
e.g.,~\cite{GripenbergLondenStaffans1990}) has a unique
solution.  The decomposition~\eqref{eq:Psi-decomp} then
gives $\theta^W$ as a function of $\Phi^J$ and $\Psi$.
The remaining freedom in $\theta^J$ given $\Phi^J$ is
clear: the map $\theta^J \mapsto \int ze^{\theta^J}\nu(dz)$
is many-to-one.
\end{proof}

\begin{remark}[Interpretation: what does the VRNC select?]
\label{rem:VRNC-selects}
The VRNC determines $\Psi$---the \emph{total effective drift} of
the memory factor $Z$ under the pricing measure.  This is the
minimal constraint needed for the stock to be priced correctly.
The remaining freedom corresponds to how jump risk is priced
\emph{within} the volatility dynamics---invisible to the stock
but visible to volatility derivatives.

Different selections of $\theta^J$ (given the same $\Psi$) produce
different $\QQ$-dynamics for volatility, hence different prices for
variance swaps and vol-of-vol products.  This is the jump analog of
the model risk gap familiar from the continuous case:
the choice of $\theta^J$ selects among
EMMs that all price the stock correctly but disagree on
vol-sensitive claims.

The VRNC reduces this infinite-dimensional ambiguity to a
structured problem: first solve the Volterra equation for $\Psi$
(the memory-consistent effective drift), then choose the
continuous--jump split according to additional criteria
(minimal entropy, variance-optimal, calibration to vol
derivatives, etc.).
\end{remark}

\begin{remark}[Connection to the defect]
\label{rem:VRNC-defect}
The defect scaling (\Cref{thm:scaling,cor:hedge-error}) and
the VRNC interact: different choices of $\theta^J$ (within the
VRNC constraint) produce different defect magnitudes in the
hedging decomposition (\Cref{thm:hedging-main}).  The constant
$C(h,\eta,p)$ in the hedging error bound depends on
$\theta^J$ through the $\QQ$-moments of volatility.  The
\emph{scaling exponents} are invariant
(\Cref{prop:defect-form-invariance}), but the constants are not.
\end{remark}


\begin{remark}[The integral form for singular kernels]
\label{rem:VRNC-singular}
For the fractional kernel $K_H(t,s) = (t-s)^{H-1/2}$ with
$H < 1/2$, $\partial_t K$ does not exist classically (it diverges
as $(t-s)^{H-3/2}$ near $s = t$).  The integral
form~\eqref{eq:VRNC-int} remains well-defined:
\[
  \int_0^t(\mu_s - r)\,ds = -\int_0^t\lambda_s\Psi_s\,ds
  - \int_0^t\beta_s\!\left[\frac{d}{ds}\int_0^s
  (s-u)^{H-1/2}\Psi_u\,du\right]ds.
\]
The inner Volterra integral $B_s^\theta = \int_0^s(s-u)^{H-1/2}
\Psi_u\,du$ is well-defined and continuous in $s$ (since
$\int_0^s(s-u)^{H-1/2}\,du = s^{H+1/2}/(H+1/2) < \infty$),
even though its derivative $\dot{B}_s^\theta$ may not exist
classically.  The integral form is therefore the primary
constraint for singular kernels.
\end{remark}


%======================================================================

\subsection{Bracket invariance and the roughness parameter}
\label{sec:bracket-inv}

We now prove that the roughness parameter $H$ is invariant
under equivalent measure change---a result of both theoretical
and practical significance.

\begin{definition}[Volterra bracket]\label{def:volterra-bracket}
For the Volterra--L\'evy process $V_T(t) = \int_0^t K(T,s)\,dL_s$,
the \emph{continuous bracket} is
\begin{equation}\label{eq:cont-bracket}
  \langle V_T^c \rangle_t
  = \sigma_L^2\int_0^t K(T,s)^2\,ds,
\end{equation}
and the \emph{jump bracket} (predictable quadratic variation
of the pure-jump part) is
\begin{equation}\label{eq:jump-bracket}
  \langle V_T^d \rangle_t
  = \int_0^t K(T,s)^2\,ds\;\cdot\;\int_\R z^2\,\nu(dz)
\end{equation}
when $\int z^2\,\nu(dz) < \infty$ (e.g., tempered stable),
and is infinite for purely stable $\nu$ with $\gamma < 2$.
The \emph{total bracket} is the predictable compensator of
$[V_T]_t = \sum_{s \leq t}(\Delta V_T(s))^2$.
\end{definition}

\begin{theorem}[Bracket invariance]\label{thm:bracket-inv}
Let $\QQ \sim \PP$ be an equivalent measure satisfying the
Girsanov and L\'evy--Girsanov conditions.  Write
$W^\QQ = W + \int_0^\cdot \theta_s^W\,ds$ for the
$\QQ$-Brownian motion and $\Ntilde^\QQ(dt,dz) =
N(dt,dz) - e^{\theta_t^J(z)}\nu(dz)\,dt$ for the $\QQ$-compensated
Poisson measure, where $\theta_t^J(z)$ may depend on time.
Then:
\begin{enumerate}[label=(\alph*)]
\item \textbf{Continuous bracket:}
  $\langle V_T^c \rangle_t^\QQ = \langle V_T^c \rangle_t^\PP
  = \sigma_L^2\int_0^t K(T,s)^2\,ds$.
  This is deterministic and identical under both measures.
\item \textbf{Jump bracket transformation:} The $\QQ$-jump
  bracket (predictable compensator of $[V_T^d]_t$) satisfies
  \begin{equation}\label{eq:jump-bracket-Q}
    \langle V_T^d \rangle_t^\QQ
    = \int_0^t K(T,s)^2
    \left(\int_\R z^2\,e^{\theta_s^J(z)}\,\nu(dz)\right)ds.
  \end{equation}
  The \emph{kernel structure} $K(T,s)^2$ within the time
  integral is unchanged; the $z^2$-moment acquires a
  time-dependent re-weighting factor
  $\int z^2 e^{\theta_s^J(z)}\nu(dz)$.
  When $\theta^J$ is time-independent, this factors as
  $(\int_0^t K(T,s)^2\,ds)\cdot(\int z^2 e^{\theta^J(z)}\nu(dz))$.
\item \textbf{Roughness invariance (continuous bracket sense):}
  The continuous bracket $\langle V_T^c\rangle_t
  = \sigma_L^2\int_0^t K(T,s)^2\,ds$ is a deterministic
  function of $t$, identical under $\PP$ and $\QQ$.
  In particular, for $K = K_H$:
  \begin{equation}\label{eq:H-invariant}
    \langle V_T^c \rangle_t
    = \frac{\sigma_L^2}{2H}
    \bigl[T^{2H} - (T-t)^{2H}\bigr]
    \qquad (\text{same under $\PP$ and $\QQ$}).
  \end{equation}
  Any quantity extracted from the continuous bracket of
  $V_T(\cdot)$ --- including the $\log$-$\log$ slope of
  $\Delta\Gamma(t,h)$ vs.\ $h$ --- is therefore
  measure-invariant.
\end{enumerate}
\end{theorem}

\begin{remark}[What is and is not proved about $H$-invariance]
\label{rem:H-scope}
\Cref{thm:bracket-inv}(c) establishes measure-invariance
of the continuous bracket of the \emph{partial process}
$V_T(t) = \int_0^t K(T,s)\,dL_s$, which is a semimartingale.
For the \emph{Volterra process} $V_t = \int_0^t K(t,s)\,dL_s$
(generally not a semimartingale), the $2H$-H\"older regularity
\[
  \E[|V_{t+h} - V_t|^2] \asymp h^{2H}
  \qquad (h \to 0)
\]
is determined by the kernel geometry near the diagonal.
This scaling is also measure-invariant when both the
continuous bracket (invariant by~(a)) and the second
jump moment $\int z^2\,\nu^\QQ(dz)$ are finite: the
latter changes by the bounded factor $e^{\theta^J(z)}$,
which rescales the prefactor but preserves the exponent.
When $\int z^2\,\nu(dz) = \infty$ (purely stable), the
second moment $\E[|V_{t+h} - V_t|^2]$ is itself infinite,
and $H$-regularity must be defined via other means (e.g.,
$p$-variation or Besov regularity).  In that setting, the
invariance of $H$ requires a separate argument, which we
do not provide here.
\end{remark}

\begin{proof}
(a) The continuous bracket $\langle V_T^c \rangle$ is the
predictable compensator of $(V_T^c)^2 = (\sigma_L\int K\,dW)^2$.
Since $K(T,\cdot)$ is deterministic, $\langle V_T^c \rangle_t
= \sigma_L^2\int_0^t K(T,s)^2\,ds$ is also deterministic.
Under the Girsanov change $W \to W^\QQ$, the continuous
bracket of a semimartingale is preserved: $\langle M^c \rangle$
is the same under $\PP$ and~$\QQ$ for any continuous
local martingale $M^c$ (this is a standard fact;
see~\cite{ContTankov2004}, Theorem~9.12).

(b) Under the L\'evy--Girsanov theorem, the compensator of
$N(dt,dz)$ changes from $\nu(dz)\,dt$ to
$e^{\theta_t^J(z)}\nu(dz)\,dt$.  The
predictable compensator of $[V_T^d]_t = \sum_{s \leq t}
K(T,s)^2(\Delta L_s)^2$ is
$\int_0^t K(T,s)^2\,(\int z^2 e^{\theta_s^J(z)}\nu(dz))\,ds$,
since the jump sizes $\Delta L_s$ retain their magnitude
under~$\QQ$ (only the compensating intensity changes).

(c) The formula~\eqref{eq:H-invariant} is an explicit computation:
$\int_0^t(T-s)^{2H-1}\,ds = [T^{2H} - (T-t)^{2H}]/(2H)$.
This is deterministic and hence identical under $\PP$
and~$\QQ$.
\end{proof}

\begin{remark}[Practical consequence]\label{rem:H-calibration}
The roughness parameter $H$ estimated from historical data
(under $\PP$) is the same $H$ that governs risk-neutral dynamics
(under~$\QQ$).  This is widely assumed in the rough volatility
literature but rarely stated as a theorem.  The proof above
shows it is a consequence of a general principle: the bracket
structure (which encodes $H$) is a geometric property of the
energy space, invariant under absolutely continuous measure
changes.  This holds for the combined Volterra--L\'evy setting,
not just the Gaussian case.
\end{remark}

\begin{corollary}[The roughness parameter is calibration-invariant]
\label{cor:roughness-greek}
The value of $H$ entering the pricing formula
$\E^\QQ[B_T^{-1}h(S_T)]$ is the same whether estimated from
historical data (under $\PP$) or from risk-neutral calibration
(under $\QQ$): the kernel $K_H$ parameterizing the energy
space is identical under both measures.

This is in contrast to parameters like the drift $\mu$ or
the vol-of-vol $\eta$, whose risk-neutral values depend on
the market price of risk $(\theta^W, \theta^J)$.  The Hurst
parameter is an energy-space parameter, not a drift parameter,
and is therefore calibration-invariant.
\end{corollary}



\subsection{Defect invariance under measure change}
\label{sec:defect-invariance}

\begin{proposition}[Structural form invariance]
\label{prop:defect-form-invariance}
The decomposition of the hedging error into $\nu$-singular and
$\nu$-regular components (\Cref{thm:hedging-main}) is
invariant under the choice of equivalent martingale
measure~$\QQ$.  Specifically:
\begin{enumerate}[label=(\alph*)]
\item The $\nu$-singular part (first-order jump exposure)
  remains $\nu$-singular under any $\QQ \sim \PP$: the
  integrand $\lambda_t z$ changes (through $\lambda_t$)
  but retains the $z$-linear structure.
\item The $\nu$-regular part (defect) remains $\nu$-regular:
  the integrand $\psi_t^{\mathrm{def}}(z)$ changes but retains
  the $z^2$-scaling for small jumps.
\item The \emph{scaling exponents} in the defect norm
  (\Cref{cor:hedge-error}) are the same for all~$\QQ$: the
  exponent $2H + 1/p$ comes from the kernel (measure-invariant
  by \Cref{thm:bracket-inv}), and the exponent
  $(2p - \gamma)^{-1/p}$ comes from the L\'evy measure's
  small-jump behavior, which is preserved up to an absolutely
  continuous Radon--Nikodym density $e^{\theta^J(z)}$.
\end{enumerate}
Only the \emph{constants} (depending on $\theta$) change
with~$\QQ$; the structural form and scaling are intrinsic.
\end{proposition}

\begin{proof}
(a) Under the L\'evy--Girsanov theorem, the $\QQ$-intensity
of jumps is $\nu^\QQ(dz) = e^{\theta^J(z)}\nu(dz)$.  The
integrand of the first-order term is $\lambda_t^\QQ\,z$,
still linear in $z$.  Since $\int|z|^p\,\nu^\QQ(dz)
= \int|z|^p e^{\theta^J(z)}\nu(dz) = \infty$ (the
singularity at $z = 0$ is unchanged by the bounded
multiplicative factor $e^{\theta^J(z)}$), the
$\nu^\QQ$-singularity persists.

(b) The defect integrand satisfies
$|\psi_t^{\mathrm{def},\QQ}(z)|
\leq C^\QQ\,(\int_t^T K^2\,ds)\,z^2$ for $|z| \leq 1$,
where $C^\QQ$ depends on $\theta$ through the
$\QQ$-moments of $\sigma$.  The $z^2$ factor is independent
of the measure change.  Since
$\int_{|z|\leq 1}|z|^{2p}\nu^\QQ(dz)
= \int_{|z|\leq 1}|z|^{2p}e^{\theta^J(z)}\nu(dz)
\leq \|e^{\theta^J}\|_\infty\int_{|z|\leq 1}
|z|^{2p}\nu(dz) < \infty$, the defect remains
$\nu^\QQ$-regular.

(c) The scaling $\int_0^T K_H(T,t)^{2p}\,dt
= T^{2pH-p+1}/(2pH-p+1)$ is deterministic and measure-free.
The small-jump moment $\int_{|z|\leq 1}|z|^{2p}\nu^\QQ(dz)$
changes by the bounded factor
$e^{\theta^J(z)} \in [e^{-\|\theta^J\|_\infty},
e^{\|\theta^J\|_\infty}]$, which does not affect the
divergence rate as $\gamma \to 2p^-$.
\end{proof}


%======================================================================

%======================================================================
\section{Scaling Law and Structural Obstructions}\label{sec:scaling}
%======================================================================

Given the pricing framework of \Cref{sec:pricing}, we now
quantify the magnitude of the unhedgeable component.  The
kernel-weighted Leibniz defect has a precise scaling in the
Hurst parameter~$H$, the horizon~$T$, and the stable
index~$\gamma$ of the L\'evy measure.  We also establish
that the $\nu$-singularity obstruction---the fact that
$L_T \notin \mathrm{im}(\widetilde{\delta}_J)$, since the
integrand $z$ has infinite $L^p(\nu)$ norm---persists
under Volterra convolution: the kernel cannot
regularize the jump-size singularity.

\subsection{Persistence of the $\nu$-singularity obstruction}

The process $L_T$ itself is
$\nu$-singular: its unique two-parameter integrand $z$ does not
belong to $\widetilde{H}_J = L^p(dt \otimes \nu \otimes dP)$
because $\int|z|^p\,\nu(dz) = \infty$ for $p < \gamma$.  We
show this obstruction persists under Volterra convolution.

\begin{proposition}[$\nu$-singularity of $V_T$]
\label{prop:VT-singular}
Let $V_T = \int_0^T K(T,s)\,dL_s$ with $K(T,\cdot) \not\equiv 0$.
Then $V_T$ is $\nu$-singular: $V_T \notin \im(\widetilde{\delta}_J)$.
\end{proposition}

\begin{proof}
By the L\'evy--It\^o decomposition,
$V_T = \sigma_L\int_0^T K(T,s)\,dW_s
+ \int_0^T\!\!\int_\R K(T,s)\,z\,\Ntilde(ds,dz)$.
The Poisson part has integrand $h(s,z) = K(T,s)\,z$.
By the uniqueness of the compensated Poisson representation
(\cite{Applebaum2009}, Theorem~4.2.5), if the Poisson part
of $V_T$ admits any representation
$\int\!\!\int \tilde{h}\,d\Ntilde$, then
$\tilde{h}(s,z) = K(T,s)\,z$ $\nu \otimes dt$-a.e.

Now compute:
\[
  \|K(T,\cdot)\,z\|_{\widetilde{H}_J}^p
  = \E\!\int_0^T\!\!\int_\R |K(T,s)|^p\,|z|^p\,\nu(dz)\,ds
  = \left(\int_0^T |K(T,s)|^p\,ds\right)
  \cdot\left(\int_\R |z|^p\,\nu(dz)\right).
\]
The first factor is finite by \Cref{ass:kernel}\ref{K1}.
The second factor is $\int_\R|z|^p\nu(dz) = \infty$
(since $p < \gamma$: the integral
$\int_{|z|\leq 1}|z|^p\,\nu_\gamma(dz)
= 2c_\gamma\int_0^1 z^{p-1-\gamma}\,dz = \infty$
when $p < \gamma$).

Thus $K(T,\cdot)z \notin \widetilde{H}_J$, and no alternative
integrand exists.  Hence
$V_T \notin \im(\widetilde{\delta}_J)$.
\end{proof}

\begin{remark}[The kernel cannot regularize the singularity]
\label{rem:kernel-cannot-regularize}
This result is not obvious \emph{a priori}: the kernel
$K(T,s)$ is bounded on compact intervals, so one might
hope it could improve the $z$-integrability.  But $K$
acts on the \emph{time} variable, while the singularity is
in the \emph{jump-size} variable $z$.  The factorization
$h(s,z) = K(T,s) \cdot z$ means the $z$-singularity is
inherited unchanged from $L_T$ to $V_T$, regardless of the
kernel's regularity in $s$.  The obstruction is an intrinsic
property of the L\'evy measure, not the temporal structure.
\end{remark}

\begin{corollary}[Two-regime decomposition persists]
\label{cor:two-regime-persists}
For any $\phi \in C_b^3(\R)$, the centered functional
$\phi(V_T) - \E[\phi(V_T)]$ decomposes into a
$\nu$-singular part (the first-order Poisson integral
$\int\!\!\int \phi'(V_T(s\!-\!))K(T,s)z\,\Ntilde(ds,dz)$)
and a $\nu$-regular part (the defect
$\Gamma_V^\phi = \widetilde{\delta}_J(g_\phi^K)$),
as in \Cref{thm:defect-decomposition}.  The kernel does not collapse
the two regimes; it merely weights the defect.
\end{corollary}


\subsection{Defect scaling for fractional kernels}

We specialize to the fractional kernel
$K_H(t,s) = (t-s)^{H-1/2}$ and derive the scaling of the
defect norm in $H$, $T$, and the L\'evy measure parameters.

\begin{lemma}[Kernel integral]\label{lem:kernel-integral}
For $K_H(T,t) = (T-t)^{H-1/2}$ and any exponent $r > 0$:
\begin{equation}\label{eq:kernel-integral}
  \int_0^T |K_H(T,t)|^r\,dt
  = \int_0^T (T-t)^{r(H-1/2)}\,dt
  = \frac{T^{r(H-1/2)+1}}{r(H-1/2)+1},
\end{equation}
provided $r(H - 1/2) > -1$, i.e.,
\begin{equation}\label{eq:H-condition}
  H > \frac{1}{2} - \frac{1}{r}.
\end{equation}
For $r = 2p$ with $p \in (1,\gamma)$, this
becomes $H > (p-1)/(2p)$.  Note that $(p-1)/(2p)$ is
\emph{positive} for $p > 1$ and increases with $p$:
at $p = 1.5$, the condition is $H > 1/6$; at $p \to 2^-$,
it approaches $H > 1/4$.  Thus $H > 0$ alone does
\emph{not} suffice; one needs $H > (p-1)/(2p)$.
\end{lemma}

\begin{assumption}[Roughness--exponent compatibility]
\label{ass:Hp}
The Hurst parameter $H$ and the working exponent $p$
satisfy $H > (p-1)/(2p)$.  This is equivalent to
$2pH - p + 1 > 0$ (convergence of the kernel integral
$\int_0^T K_H(T,t)^{2p}\,dt$).  Note:
\begin{itemize}
\item When $H \geq 1/2$: the condition is vacuous
  (satisfied for all $p \geq 1$).
\item When $H < 1/2$: the condition requires
  $p < 1/(1 - 2H)$.  For example, $H = 0.1$ requires
  $p < 1.25$; $H = 0.25$ requires $p < 2$.
\end{itemize}
\end{assumption}

\begin{theorem}[Defect scaling law]\label{thm:scaling}
Let $K = K_H$ (fractional kernel), $\nu$ satisfying
\Cref{ass:nu} with local stable index $\gamma \in (1,2)$,
and $f \in C_b^3(\R)$.  Fix $p \in (1,\gamma)$ satisfying
\Cref{ass:Hp} (i.e., $H > (p-1)/(2p)$), and assume the
tail condition~\ref{L2} holds with $q \geq p$.  Then:
\begin{equation}\label{eq:scaling-law}
  \|\Gamma_V^f\|_{L^p}^p
  \leq C_p^p\,\|g_f^{K_H}\|_{\widetilde{H}_J}^p
  \leq C(f,p)\;\frac{T^{2pH - p + 1}}{2pH - p + 1}
  \;\cdot\;\frac{2c_\gamma}{2p - \gamma}
  \;+\; R(f,T,p),
\end{equation}
where:
\begin{itemize}
\item The first inequality is the L\'evy BDG upper bound
  (with $C_p$ the BDG constant from~\cite{Dirksen2015}).
\item $C(f,p) = C_p^p \cdot \|f''\|_\infty^p / 2^p$
  is the small-jump coefficient.
\item $R(f,T,p)$ is the large-jump remainder:
  $R = C'\,T\!\int_{|z|>1}(1+|z|^p)\nu(dz)$, which is
  finite under~\ref{L2} with $q \geq p$.
\end{itemize}
In particular, the defect $L^p$-norm (to leading order) scales as:
\begin{equation}\label{eq:scaling-summary}
  \|\Gamma_V^f\|_{L^p}
  \sim C \cdot T^{2H - 1 + 1/p}
  \cdot \left(\frac{c_\gamma}{2p - \gamma}\right)^{1/p}
  \qquad \text{as } T \to \infty \text{ or } \gamma \to 2p^-.
\end{equation}
\end{theorem}

\begin{proof}
By the two-parameter BDG inequality (\cite{Dirksen2015},
Corollary~3.3),
$\|\widetilde{\delta}_J(g)\|_{L^p} \leq
C_p\,\|g\|_{\widetilde{H}_J}$.
The norm bound~\eqref{eq:gfK-bound} with the kernel integral
(\Cref{lem:kernel-integral}, $r = 2p$):
\[
  \int_0^T |K_H(T,t)|^{2p}\,dt
  = \frac{T^{2p(H-1/2)+1}}{2p(H-1/2)+1}
  = \frac{T^{2pH - p + 1}}{2pH - p + 1},
\]
and the jump moment:
$\int_{|z|\leq 1}|z|^{2p}\,\nu_\gamma(dz)
= 2c_\gamma\int_0^1 z^{2p-1-\gamma}\,dz
= 2c_\gamma/(2p - \gamma)$,
yield the bound~\eqref{eq:scaling-law}.

For the large-$T$ or near-critical ($\gamma \to 2p$) behavior,
the large-jump remainder is subdominant, giving~\eqref{eq:scaling-summary}.
\end{proof}


\begin{corollary}[Roughness--jump interaction]
\label{cor:interaction}
The defect norm is controlled by the product
\[
  \|\Gamma_V^f\|_{L^p} \lesssim
  \underbrace{T^{2H-1+1/p}}_{\text{roughness contribution}}
  \times
  \underbrace{(2p - \gamma)^{-1/p}}_{\text{jump contribution}}.
\]
The interaction is multiplicative:
\begin{enumerate}[label=(\alph*)]
\item \textbf{Pure rough volatility} ($\nu = 0$, $\sigma > 0$):
  $\Gamma_V^f = 0$.  The continuous energy space suffices,
  and perfect hedging is achievable (recovering the result
  of~\cite{ElEuchRosenbaum2019}).
\item \textbf{Markov limit} ($H = 1/2$):
  $T^{2H-1+1/p} = T^{1/p}$.  The defect reduces to the
  standard L\'evy Leibniz defect of~\cite{Fontes2026d},
  without kernel amplification.
\item \textbf{Rough-jump regime} ($H < 1/2$, $\gamma < 2$):
  the exponent $2H - 1 + 1/p < 1/p$, so the defect
  grows more slowly with $T$ than in the Markov case.

  \textbf{However}, the kernel singularity $K_H(T,t) =
  (T-t)^{H-1/2}$ amplifies \emph{recent} jumps: for
  $t$ near $T$, $K_H(T,t) \gg 1$, so the integrand
  $g_f^K(t,z) \approx \frac{1}{2}f''(T-t)^{2(H-1/2)}z^2$
  is large.  The defect norm concentrates on a
  neighborhood of $T$ of width $\sim \epsilon$, with
  the local contribution scaling as
  $\epsilon^{2p(H-1/2)+1}$---which is finite but large
  when $H \ll 1/2$.
\end{enumerate}
\end{corollary}

\begin{remark}[The interaction is not additive]
\label{rem:not-additive}
In a model with roughness $H$ and jump index $\gamma$, the
unhedgeable residual is \emph{not}
``rough-vol-error $+$ jump-error.''  The kernel $K_H$ enters the
defect integrand $g_f^K(t,z) = f(V_T(t\!-\!) + K_H(T,t)z)
- f(V_T(t\!-\!)) - f'(V_T(t\!-\!))K_H(T,t)z$ multiplicatively:
larger kernels produce larger Taylor remainders for each jump.
This multiplicative structure is a genuinely new phenomenon
visible only in the combined Volterra--L\'evy setting.
\end{remark}


\subsection{Sharpness: the variance swap payoff}

The upper bound in \Cref{thm:scaling} arises from two
inequalities: the Taylor bound $|g_f^K| \leq
\frac{1}{2}\|f''\|_\infty K^2 z^2$ for small jumps, and
the BDG inequality $\|\widetilde{\delta}_J(g)\|_{L^p}
\leq C_p\|g\|_{\widetilde{H}_J}$.  We now show
that the scaling exponents are sharp by exhibiting a payoff
for which the Taylor bound is exact.

\begin{proposition}[Sharpness of scaling: quadratic payoff]
\label{thm:sharpness}
Let $f(x) = x^2$, $K = K_H$, and assume:
\begin{enumerate}[label=(\roman*)]
\item $H > (p-1)/(2p)$ (\Cref{ass:Hp});
\item $\nu_{2p} := \int_\R |z|^{2p}\,\nu(dz) < \infty$
  (i.e., the tail condition~\ref{L2} holds with $q \geq 2p$;
  this is \emph{not} automatic under purely stable tails,
  but holds under tempered stable tails).
\end{enumerate}
Note that $f \notin C_b^3(\R)$; however, the defect
integrand $g_f^K$ is deterministic and has polynomial
growth in $z$, so the following identities hold by
direct computation (not by invoking \Cref{thm:VL-ito},
which assumes $f \in C_b^3$):
\begin{enumerate}[label=(\alph*)]
\item The defect integrand is exact:
  \begin{equation}\label{eq:gfK-exact}
    g_f^K(t,z) = K_H(T,t)^2\,z^2
    \qquad \text{(deterministic)}.
  \end{equation}
\item The $\widetilde{H}_J$-norm is
  \begin{equation}\label{eq:gfK-exact-norm}
    \|g_f^K\|_{\widetilde{H}_J}^p
    = \int_0^T |K_H(T,t)|^{2p}\,dt
    \;\cdot\;\int_\R |z|^{2p}\,\nu(dz)
    = \frac{T^{2pH - p + 1}}{2pH - p + 1}
    \;\cdot\;\nu_{2p}.
  \end{equation}
  Under tempered stable tails, $\nu_{2p}$ has both small-jump
  and large-jump contributions:
  $\nu_{2p} = \int_{|z|\leq 1}|z|^{2p}\nu(dz)
  + \int_{|z|>1}|z|^{2p}\nu(dz)$;
  the first equals $2c_\gamma/(2p-\gamma)$ and the second
  is finite by the exponential tail decay.  The leading
  singularity as $\gamma \to 2p^-$ is from the small-jump
  part.
\item The $L^2$ isometry (under the additional condition
  $\nu_4 := \int_\R z^4\,\nu(dz) < \infty$ and $H > 1/4$):
  \begin{equation}\label{eq:Gamma-L2}
    \|\Gamma_V^f\|_{L^2}^2
    = \int_0^T K_H(T,t)^4\,dt\;\cdot\;\nu_4
    = \frac{T^{4H-1}}{4H - 1}\;\cdot\;\nu_4.
  \end{equation}
  This is an \emph{equality} (It\^o isometry for
  deterministic integrands), not merely an inequality.
\end{enumerate}
\end{proposition}

\begin{proof}
(a) $f(x) = x^2$ gives $f(x + y) - f(x) - 2xy = y^2$.
Set $x = V_T(t\!-\!)$, $y = K_H(T,t)z$:
$g_f^K(t,z) = K_H(T,t)^2 z^2$.  This is deterministic
(no dependence on $V_T(t\!-\!)$).  The identity is exact for
all $z$, with no Taylor remainder.  That $f \notin C_b^3$ is
immaterial: the algebraic identity $(x+y)^2 - x^2 - 2xy = y^2$
holds pointwise.

(b) Since $g_f^K$ is deterministic, the expectation in
the $\widetilde{H}_J$-norm is trivial:
$\|g_f^K\|_{\widetilde{H}_J}^p
= \int_0^T |K_H(T,t)|^{2p}\,dt \cdot \int_\R|z|^{2p}\nu(dz)$.
The time integral converges by \Cref{ass:Hp}; the jump integral
converges by assumption~(ii).

(c) For deterministic integrands $h \in L^2(dt \otimes \nu)$,
the compensated Poisson integral satisfies the isometry
$\E[|\widetilde{\delta}_J(h)|^2]
= \int_0^T\!\int_\R h(t,z)^2\,\nu(dz)\,dt$
(\cite{Applebaum2009}, Theorem~4.2.3).  Setting
$h(t,z) = K_H(T,t)^2 z^2$ and computing:
$\int_0^T K_H(T,t)^4\,dt = T^{4H-1}/(4H-1)$ (by
\Cref{lem:kernel-integral} with $r = 4$, requiring $H > 1/4$),
and $\int_\R z^4\,\nu(dz) = \nu_4 < \infty$ by assumption.
\end{proof}

\begin{remark}[Why $f(x)=x^2$ is legitimate despite $f \notin C_b^3$]
\label{rem:quadratic-legitimate}
The Volterra--L\'evy It\^o formula (\Cref{thm:VL-ito}) is
stated for $f \in C_b^3$.  However, for $f(x) = x^2$, the
defect integrand $g_f^K(t,z) = K(T,t)^2 z^2$ is the exact
finite difference, not a Taylor approximation.  The It\^o
formula for the semimartingale $V_T(t)$ applied to $x^2$
is valid under the moment condition
$\E[\int_0^T |K(T,t)|^2|\Delta L_t|^2] < \infty$,
which holds when $\nu_2 < \infty$ (tempered stable) or
by stopping time arguments for purely stable drivers.
The bounded-smooth assumption in \Cref{thm:VL-ito} is
used for the general Taylor-remainder bound; for the
specific quadratic case, it is unnecessary.
\end{remark}

\begin{remark}[Significance for variance swaps]
\label{rem:var-swap-sharpness}
\Cref{thm:sharpness} shows that the
upper bound in \Cref{thm:scaling} is \emph{achieved} (up
to BDG constants) for the quadratic payoff class.  The
scaling exponent $2pH - p + 1$ in $T$ is therefore tight:
it is not an artifact of the Taylor approximation.
The dependence on $\nu_{2p}$ is also exact, though
the decomposition into small-jump and large-jump
contributions depends on the tail behavior of $\nu$.
\end{remark}

\begin{remark}[BDG sharpness]\label{rem:BDG-sharp}
For $p \neq 2$, the BDG inequality
for Poisson integrals has a two-sided form
$c_p\|h\|_{\widetilde{H}_J} \leq
\|\widetilde{\delta}_J(h)\|_{L^p} \leq
C_p\|h\|_{\widetilde{H}_J}$ when $h$ is deterministic
(see~\cite{Dirksen2015}, Corollary~3.3).  Under the
conditions of \Cref{thm:sharpness}, this gives:
\begin{equation}\label{eq:two-sided}
  c_p\;\frac{T^{2H - 1 + 1/p}}{(2pH-p+1)^{1/p}}
  \;\cdot\;\nu_{2p}^{1/p}
  \leq \|\Gamma_V^{x^2}\|_{L^p}
  \leq C_p\;\frac{T^{2H - 1 + 1/p}}{(2pH-p+1)^{1/p}}
  \;\cdot\;\nu_{2p}^{1/p}.
\end{equation}
\end{remark}


\section{Application to Hedging}\label{sec:hedging}
%======================================================================

This section makes the connection between the Volterra--L\'evy
It\^o formula and derivatives hedging precise.  The core
difficulty is the \emph{composition chain}: a jump of size $z$ at
time $t$ in the L\'evy driver propagates through the kernel $K$
into the volatility process, then through the exponential
vol-of-vol map into instantaneous volatility, then through the
stock price dynamics into the terminal payoff.  Each nonlinear
link in this chain generates a defect term; the total hedging
residual is their aggregate.


\subsection{Rough-jump stochastic volatility model}

\begin{assumption}[Market model]\label{ass:market}
The filtration $(\F_t)_{t \in [0,T]}$ is generated by three
independent sources of noise:
\begin{enumerate}[label=(N\arabic*)]
\item\label{N1} A Brownian motion $W^\perp = (W_t^\perp)$
  (``idiosyncratic stock noise'');
\item\label{N2} A Brownian motion $W = (W_t)$
  (``continuous volatility driver'');
\item\label{N3} A compensated Poisson measure
  $\Ntilde(dt,dz)$ on $[0,T] \times \R$ with intensity
  $\nu$ satisfying \Cref{ass:nu}
  (``jump volatility driver'').
\end{enumerate}
The three sources are mutually independent under the physical
measure $\PP$.
\end{assumption}

The asset price and volatility dynamics are:
\begin{align}
  \frac{dS_t}{S_t} &= \mu_t\,dt + \sigma_t\,dW_t^S,
  \label{eq:S-dynamics}\\
  \log \sigma_t &= \log \sigma_0 + \eta\,V_t
  - \tfrac{1}{2}\eta^2\,\E[V_t^2],
  \label{eq:vol-dynamics}\\
  V_t &= \int_0^t K(t,s)\,dL_s, \qquad
  L_t = \sigma_L\,W_t + J_t,
  \label{eq:V-dynamics}
\end{align}
where:
\begin{enumerate}[label=(\roman*)]
\item $W^S = \rho\,W + \sqrt{1 - \rho^2}\,W^\perp$ is the
  stock-driving Brownian motion, with leverage parameter
  $\rho \in (-1,1)$.
\item $J_t = \int_0^t\!\!\int_\R z\,\Ntilde(ds,dz)$ is the
  pure-jump component of the L\'evy driver, $\sigma_L \geq 0$.
\item The kernel $K$ satisfies \Cref{ass:kernel}.
\item The centering $-\frac{1}{2}\eta^2\E[V_t^2]$ ensures
  $\E[\sigma_t^2] = \sigma_0^2$ (martingale correction for the
  exponential).
\end{enumerate}

\begin{remark}[Independence structure]\label{rem:independence}
The stock noise $W^S$ is a linear combination of the two Brownian
components.  Equivalently, $W = \rho\,W^S
+ \sqrt{1-\rho^2}\,W^\perp_\star$ for an independent Brownian
$W^\perp_\star$, but the $(W^\perp, W, \Ntilde)$ basis is more
natural for the operator decomposition: each component generates
a separate energy space, and the operator derivatives split
accordingly (\Cref{prop:D-split}).
\end{remark}


\subsection{Martingale representation under three drivers}

Under an equivalent pricing measure $\QQ$ (see
\Cref{rem:Q-existence} below), the discounted stock
$\tilde{S}_t := e^{-rt}S_t$ is a $\QQ$-martingale.  Write
$W^{\perp,\QQ}$, $W^\QQ$, and $\Ntilde^\QQ$ for the
$\QQ$-Brownian motions and compensated Poisson measure.
By Girsanov's theorem (applied to the Brownian components)
and the L\'evy--Girsanov theorem (for $\Ntilde$;
see~\cite{ContTankov2004}, Theorem~9.12), the three
sources remain independent under~$\QQ$.

\begin{proposition}[Three-driver MRT]\label{prop:MRT}
Under \Cref{ass:market}, every $F \in L^2(\F_T, \QQ)$ admits
the unique martingale representation
\begin{equation}\label{eq:MRT}
  F = \E^\QQ[F]
  + \int_0^T \phi_t^\perp\,dW_t^{\perp,\QQ}
  + \int_0^T \phi_t^W\,dW_t^\QQ
  + \int_0^T\!\!\int_\R
  \psi_t(z)\,\Ntilde^\QQ(dt,dz),
\end{equation}
where $\phi^\perp, \phi^W$ are predictable processes in
$L^2(\sigma^2\,dt \otimes d\QQ)$ and $\psi$ is a predictable
random field in $L^2(dt \otimes \nu^\QQ(dz) \otimes d\QQ)$.
\end{proposition}

\begin{proof}
The filtration $\F_t$ is generated by the
three independent processes
$(W^{\perp,\QQ}, W^\QQ, \Ntilde^\QQ)$ under~$\QQ$.  The joint
predictable representation property follows from the PRP of
Brownian motion (for each $W^{\perp,\QQ}$ and $W^\QQ$) and the
martingale representation for Poisson random measures
(\cite{Applebaum2009}, Theorem~5.3.5).  Independence ensures
uniqueness of the three-component decomposition.
\end{proof}

\begin{remark}[Existence of $\QQ$]\label{rem:Q-existence}
The model~\eqref{eq:S-dynamics}--\eqref{eq:V-dynamics} is
generically \emph{incomplete}: three noise sources drive
one traded asset.  Existence of an EMM $\QQ$ requires the
market price of risk $\theta = (\theta^\perp, \theta^W,
\theta^J(\cdot))$ to satisfy the integrability conditions
of the Girsanov and L\'evy--Girsanov theorems.  For the
Brownian components, this is the standard Novikov condition.
For the jump component, the Esscher-type condition
$\int_\R(e^{\theta^J(z)} - 1 - \theta^J(z))\,\nu(dz) < \infty$
suffices.  The Volterra risk-neutral constraint
(\Cref{thm:VRNC-VL}) provides additional structural conditions
when memory factors are present.  The hedging analysis below
holds for any fixed $\QQ$ satisfying these conditions; the
\emph{choice} of $\QQ$ among valid EMMs affects the hedging
integrands $\phi^\perp, \phi^W$ but not the structural form of
the defect or its scaling.
\end{remark}


\subsection{From MRT to traded instruments}

The representation~\eqref{eq:MRT} decomposes $h(S_T)$ into three
orthogonal components, each associated to a noise source.  We now
show how the first two are absorbed by stock and variance swap
trading, and how the third---the jump component---carries the
kernel-weighted Leibniz defect.

\subsubsection{Stock hedge (absorbing $dW^{\perp,\QQ}$ and
part of $dW^\QQ$)}

Under $\QQ$, the stock dynamics are
\begin{equation}\label{eq:SQ}
  dS_t = S_t\bigl[r\,dt + \sigma_t\,(\rho\,dW_t^\QQ
  + \sqrt{1-\rho^2}\,dW_t^{\perp,\QQ})\bigr].
\end{equation}
Define the \emph{stock delta process}
$\Delta_t := \frac{1}{S_t\sigma_t}
\bigl(\sqrt{1-\rho^2}\;\phi_t^\perp + \rho\;\phi_t^W\bigr)$.
Then
\begin{equation}\label{eq:stock-hedge-absorption}
  \Delta_t\,dS_t = \Delta_t\,S_t\,\sigma_t\,dW_t^{S,\QQ} + \ldots
  = (\sqrt{1-\rho^2}\;\phi_t^\perp + \rho\;\phi_t^W)\,
  (\rho\,dW_t^\QQ + \sqrt{1-\rho^2}\,dW_t^{\perp,\QQ})
  + \ldots
\end{equation}
Expanding and using the orthogonality $dW^\QQ \cdot dW^{\perp,\QQ} = 0$:
\begin{equation}\label{eq:stock-absorbs}
  \Delta_t\,dS_t \;=\;
  (1-\rho^2)\,\phi_t^\perp\,dW_t^{\perp,\QQ}
  + \rho^2\,\phi_t^W\,dW_t^\QQ
  + \rho\sqrt{1-\rho^2}\,
  (\phi_t^\perp + \phi_t^W)\,dW_t^\QQ
  + \ldots
\end{equation}

We can simplify this by working in the rotated basis.  The stock
hedge absorbs \emph{all} of the $dW^{S,\QQ}$-exposure:
\begin{equation}\label{eq:stock-full}
  \Delta_t\,dS_t = \phi_t^S\,dW_t^{S,\QQ} + (\text{drift}),
  \quad \text{where }
  \phi_t^S := \sqrt{1-\rho^2}\;\phi_t^\perp + \rho\;\phi_t^W.
\end{equation}
What remains after stock hedging is the component of $\phi^W$
orthogonal to the stock:
\begin{equation}\label{eq:residual-after-stock}
  \phi_t^{W,\perp} := \phi_t^W
  - \frac{\rho}{\sqrt{1-\rho^2}}\,\phi_t^\perp
  \cdot \frac{\sqrt{1-\rho^2}}{\sigma_t S_t}
  \cdot \sigma_t S_t\,\rho,
\end{equation}
which simplifies: the unhedged continuous exposure after delta
hedging is
\begin{equation}\label{eq:after-delta}
  R_t^{\mathrm{cont}} := \phi_t^{W,\perp}\,dW_t^{\perp_\star,\QQ},
\end{equation}
where $W^{\perp_\star}$ is the Brownian motion orthogonal to
$W^S$ in the $(W,W^\perp)$-plane.  This is the \emph{pure
volatility exposure}---the component of the claim's value
driven by volatility innovations that are uncorrelated with
the stock.

Rather than tracking the rotation explicitly, we state the
result in coordinate-free form:

\begin{lemma}[Residual after delta hedging]\label{lem:after-delta}
Let $F = h(S_T)$ with MRT~\eqref{eq:MRT}.  After continuously
delta-hedging with the stock, the hedging error is
\begin{equation}\label{eq:delta-error}
  h(S_T) - \E^\QQ[h(S_T)]
  - \int_0^T \Delta_t\,dS_t
  = \int_0^T \phi_t^{W,\perp}\,dW_t^{\perp_\star,\QQ}
  + \int_0^T\!\!\int_\R
  \psi_t(z)\,\Ntilde^\QQ(dt,dz),
\end{equation}
where $\phi^{W,\perp} =
(1-\rho^2)^{-1/2}(\phi^W - \rho\,\phi^\perp)$ is the
orthogonal volatility integrand.  In particular, the jump
component is entirely unaffected by stock trading.
\end{lemma}

\begin{proof}
The stock $S$ has zero loading on $\Ntilde$ and a specific
linear combination of loadings on $(W^\perp, W)$.  Delta
hedging projects onto the $W^S$-direction; the orthogonal
complement in the Brownian plane and the entire Poisson
component remain.  This follows from the independence of
$(W^{\perp_\star,\QQ}, \Ntilde^\QQ)$ and $W^{S,\QQ}$
under~$\QQ$.
\end{proof}


\subsubsection{Variance swap hedge (absorbing $dW^{\perp_\star,\QQ}$)}

The pure volatility exposure $\phi^{W,\perp}\,dW^{\perp_\star,\QQ}$
can be hedged by trading a variance swap (or more generally, any
derivative whose value depends on the volatility path).

Let $\Xi_t := \E^\QQ[\int_0^T \sigma_s^2\,ds \mid \F_t]$ denote
the fair-value process of a variance swap struck at inception.
Since $\sigma_t$ depends on the continuous volatility driver $W$
through~\eqref{eq:vol-dynamics}, $\Xi_t$ has a nonzero loading
on $dW^{\perp_\star,\QQ}$.

\begin{lemma}[Variance swap exposure to $W^{\perp_\star}$]
\label{lem:xi-loading}
Under the exponential specification~\eqref{eq:vol-dynamics}, the
dynamics of $\Xi$ contain a term
\begin{equation}\label{eq:xi-continuous}
  d\Xi_t^{\mathrm{cont}} = \xi_t^W\,dW_t^{\perp_\star,\QQ}
  + \xi_t^S\,dW_t^{S,\QQ} + \ldots
\end{equation}
with $\xi_t^W \neq 0$ a.s.\ whenever $\eta\sigma_L > 0$
(i.e., whenever the model has stochastic volatility with a
continuous component).
\end{lemma}

\begin{proof}[Proof sketch]
By the tower property,
$\Xi_t = \int_0^t \sigma_s^2\,ds
+ \E^\QQ[\int_t^T \sigma_s^2\,ds \mid \F_t]$.
The first term is adapted and does not contribute to
$d\Xi^{\mathrm{mart}}$.
For the conditional expectation, the martingale
representation theorem provides integrands against each
of the three noise sources.  Under the exponential
specification, $\sigma_s = \sigma_0\exp(\eta V_s
- \frac{1}{2}\eta^2\E[V_s^2])$ and $V_s$ loads on $W$
through the kernel $K$, so the martingale part of $\Xi$
inherits a $dW^{\perp_\star,\QQ}$ loading proportional to
$\eta\sigma_L\int_t^T\!\int_t^s K(s,u)^2\,du\;
\sigma_u^2\,ds$, which is a.s.\ positive.
A rigorous derivation requires the explicit martingale
representation of $\sigma_s^2$ (and hence of $\Xi$) with
controlled $L^2$-integrability of the resulting integrands;
we assume these regularity conditions are satisfied under
\Cref{ass:exp-moment}.
\end{proof}

By trading $\xi^W$-units of the variance swap per unit of
orthogonal vol exposure, we absorb the
$\phi^{W,\perp}\,dW^{\perp_\star,\QQ}$ component.  Define
$\Delta_t^\Xi :=
\phi_t^{W,\perp}\,/\,\xi_t^W$.

\begin{proposition}[Residual after stock + variance swap hedging]
\label{prop:residual-SV}
After continuously delta-hedging with the stock and vega-hedging
with the variance swap:
\begin{equation}\label{eq:SV-residual}
  h(S_T) - \E^\QQ[h(S_T)]
  - \int_0^T \Delta_t\,dS_t
  - \int_0^T \Delta_t^\Xi\,d\Xi_t
  = \int_0^T\!\!\int_\R \psi_t(z)\,\Ntilde^\QQ(dt,dz)
  + \varepsilon_T^{\Xi},
\end{equation}
where $\varepsilon_T^\Xi$ accounts for the jump component of the
variance swap dynamics (which we absorb into the Poisson integral
below).  The continuous hedging error vanishes: stock $+$ variance
swap absorbs all of $\phi^\perp\,dW^{\perp,\QQ}$ and
$\phi^W\,dW^\QQ$.
\end{proposition}

\begin{proof}
The stock trade absorbs $\phi^S\,dW^{S,\QQ}$
(\Cref{lem:after-delta}).  The variance swap trade absorbs
$\phi^{W,\perp}\,dW^{\perp_\star,\QQ}$ by construction.
Together, $\phi^S\,dW^{S,\QQ}
+ \phi^{W,\perp}\,dW^{\perp_\star,\QQ}$ spans the
full Brownian part of the MRT~\eqref{eq:MRT}.

The variance swap $\Xi$ also has a jump component: when $V$
jumps, $\sigma$ jumps, and $\Xi$ (as a conditional expectation)
adjusts.  The term $\varepsilon_T^\Xi$ collects the jump-on-jump
cross terms from the variance swap trade.  Since these are also
driven by $\Ntilde^\QQ$, they are absorbed into the total
Poisson integrand $\psi_t^{\mathrm{total}}(z)
:= \psi_t(z) - \Delta_t^\Xi\,\Delta\Xi_t(z)$, where
$\Delta\Xi_t(z)$ is the jump in $\Xi$ caused by a jump of
size $z$ at time~$t$.
\end{proof}


\subsection{The jump residual: composition and defect}

What remains after stock and variance swap hedging is a
compensated Poisson integral.  We now identify its structure
by tracing the composition chain explicitly.

\subsubsection{The composition chain}

A jump of size $z$ in the Poisson measure at time $t < T$
propagates as follows:
\begin{equation}\label{eq:composition-chain}
  z \;\xrightarrow{\;\;K\;\;}
  \Delta V_s = K(s,t)\,z \;\text{ for all }s \geq t
  \;\xrightarrow{\;\;\exp\;\;}
  \sigma_s \mapsto \sigma_s\,e^{\eta K(s,t)z}
  \;\xrightarrow{\;\;\int\sigma\,dW^S\;\;}
  \Delta(\log S_T) \neq 0
  \;\xrightarrow{\;\;h\;\;}
  \Delta h(S_T).
\end{equation}

Crucially, a single jump at time $t$ affects the volatility
at all future times $s \in [t,T]$ through the kernel weight
$K(s,t)$, and hence affects $S_T$ through the integrated
volatility $\int_t^T \sigma_s^2\,ds$ and the stochastic
integral $\int_t^T \sigma_s\,dW_s^S$.  This is the
\emph{memory amplification}: the kernel converts a pointwise
jump into a path-level shift.


\subsubsection{The jump integrand of the claim}

The following lemma establishes the regularity of the
scenario map that the jump integrand decomposition requires.

\begin{lemma}[Regularity of the scenario map]\label{lem:scenario-reg}
Under \Cref{ass:market,ass:nu,ass:kernel,ass:exp-moment},
let $h \in C_b^3(\R)$.  Then the scenario map
$z \mapsto S_T^{t,z}$ (the terminal stock price when
$\Ntilde$ has an atom of size $z$ at time~$t$) is
twice differentiable at $z = 0$ in conditional
$L^{2p}(\QQ \mid \F_{t-})$, and the derivative
$\lambda_t := \partial_z \psi_t(0)$ is predictable.
\end{lemma}

\begin{proof}
The composition chain~\eqref{eq:composition-chain} has
three links; we verify differentiability and moment
control at each.

\textbf{Volatility link.}  By the exponential
specification~\eqref{eq:vol-dynamics},
$\sigma_s^{t,z} = \sigma_s\,e^{\eta K(s,t)z}$.
For every fixed $s \geq t$, the map $z \mapsto
\sigma_s^{t,z}$ is $C^\infty$ in $z$, with
$\partial_z \sigma_s^{t,z} = \eta K(s,t)\,\sigma_s^{t,z}$.
Under~\ref{E1}, $\E[\sigma_s^{t,z\,;\,r}] < \infty$ for
all $r > 0$ and $|z| \leq 1$, so the derivatives
are uniformly $L^r$-bounded in $z$ on $[-1,1]$.

\textbf{Stock link.}  By~\eqref{eq:logS-shift},
$\log S_T^{t,z} - \log S_T^{t,0}
= \int_t^T \sigma_s(e^{\eta K(s,t)z}-1)\,dW_s^S
- \frac{1}{2}\int_t^T \sigma_s^2(e^{2\eta K(s,t)z}-1)\,ds$.
The BDG inequality gives, for $|z| \leq 1$:
\[
  \E^{\QQ}\!\left[\left.
  |S_T^{t,z} - S_T^{t,0}|^{2p}\;\right|\;\F_{t-}\right]
  \leq C\,S_T^{2p}\,\eta^{2p}\,z^{2p}\,
  \E^{\QQ}\!\left[\left.
  \left(\int_t^T K(s,t)^2\sigma_s^2\,ds\right)^p
  \;\right|\;\F_{t-}\right],
\]
which is finite under~\ref{E2}.  Differentiability
of $z \mapsto S_T^{t,z}$ in conditional $L^{2p}$
follows by dominated convergence applied to the
difference quotients, using the uniform $L^{2p}$
bound above.

\textbf{Payoff link.}  Since $h \in C_b^3$,
$z \mapsto h(S_T^{t,z})$ inherits the differentiability
of $z \mapsto S_T^{t,z}$ by the chain rule, with
$|\partial_z^2 h(S_T^{t,z})| \leq \|h''\|_\infty
\cdot |\partial_z S_T^{t,z}|^2
+ \|h'\|_\infty \cdot |\partial_z^2 S_T^{t,z}|$,
both in $L^p$ by the stock-link estimates.

\textbf{Predictability.}  The derivative
$\lambda_t = \partial_z \psi_t(0)$ is a conditional
expectation of $\F_T$-measurable random variables given
$\F_{t-}$, hence $\F_{t-}$-measurable and predictable.
\end{proof}

\begin{proposition}[Jump integrand decomposition]\label{thm:jump-decomp}
Under \Cref{ass:market,ass:nu,ass:kernel,ass:exp-moment},
let $h \in C_b^3(\R)$.
Then the MRT jump integrand $\psi_t^{\mathrm{total}}(z)$
in~\eqref{eq:SV-residual} (after absorbing the variance swap
jump correction) decomposes as
\begin{equation}\label{eq:psi-decomp}
  \psi_t^{\mathrm{total}}(z) = \psi_t^{\mathrm{lin}}(z)
  + \psi_t^{\mathrm{def}}(z),
\end{equation}
where:
\begin{enumerate}[label=(\alph*)]
\item The \emph{linear part} $\psi_t^{\mathrm{lin}}(z)$ is linear
  in $z$ and is the first-order sensitivity of the (delta- and
  vega-hedged) claim value to a jump of size $z$ at time $t$:
  \begin{equation}\label{eq:psi-lin}
    \psi_t^{\mathrm{lin}}(z) = \lambda_t \cdot z,
  \end{equation}
  where $\lambda_t$ is a predictable process depending on the
  current state $(\sigma_t, S_t, \Xi_t)$ and the kernel
  $K(\cdot, t)$.

\item The \emph{defect part} $\psi_t^{\mathrm{def}}(z)$ is the
  aggregate nonlinear response---the kernel-weighted Leibniz
  defect of the composition:
  \begin{equation}\label{eq:psi-def}
    \psi_t^{\mathrm{def}}(z)
    = \E^\QQ\!\left[\left.
    B_T^{-1}h(S_T^{t,z}) - B_T^{-1}h(S_T^{t,0})
    - \lambda_t\,z\;\right|\;\F_{t-}\right],
  \end{equation}
  where $S_T^{t,z}$ denotes the terminal stock price in a
  scenario where $\Ntilde$ has an atom of size $z$ at time~$t$,
  and $S_T^{t,0}$ is the same scenario without the atom.
\end{enumerate}
\end{proposition}

\begin{proof}
By \Cref{lem:scenario-reg}, $z \mapsto S_T^{t,z}$ is
twice differentiable in conditional $L^{2p}$, so the
first-order extraction below is justified.

The MRT~\eqref{eq:MRT} is constructive: the Poisson integrand
$\psi_t(z)$ is the predictable projection of the finite
difference in the value process caused by a jump:
\begin{equation}\label{eq:psi-FD}
  \psi_t(z) = \E^\QQ\!\left[\left.
  B_T^{-1}h(S_T^{t,z}) - B_T^{-1}h(S_T^{t,0})
  \;\right|\;\F_{t-}\right].
\end{equation}
This is the standard representation of the Poisson MRT integrand
(see~\cite{Applebaum2009}, Theorem~5.3.5;
\cite{ContTankov2004}, Proposition~14.23).

The total integrand $\psi^{\mathrm{total}}$ in~\eqref{eq:SV-residual}
subtracts the variance swap jump contribution, but retains the
same finite-difference structure with the hedged claim value
replacing the unhedged one.

Define the linear part by extracting the first-order Taylor
expansion in $z$:
\[
  \psi_t^{\mathrm{lin}}(z) := \left.\frac{\partial}{\partial z}
  \psi_t(z)\right|_{z=0} \cdot z
  = \lambda_t\,z,
\]
where $\lambda_t = \partial_z \psi_t(0)$ is the
\emph{jump delta}---the marginal sensitivity of the
hedged claim to an infinitesimal jump.  Then
$\psi_t^{\mathrm{def}}(z)
:= \psi_t^{\mathrm{total}}(z) - \psi_t^{\mathrm{lin}}(z)$
is the remainder.
\end{proof}


\subsubsection{Integrability: the two-regime structure}

\begin{proposition}[Integrability dichotomy]\label{prop:integrability}
Under the stable-type \Cref{ass:nu} with $\gamma \in (1,2)$:
\begin{enumerate}[label=(\alph*)]
\item $\psi^{\mathrm{lin}}(t,z) = \lambda_t\,z$ is
  \emph{$\nu$-singular}: it does not belong to
  $\widetilde{H}_J = L^p(dt \otimes \nu \otimes d\QQ)$.

\item $\psi^{\mathrm{def}}(t,z)$ is
  \emph{$\nu$-regular}: it belongs to $\widetilde{H}_J$.
\end{enumerate}
\end{proposition}

\begin{proof}
(a) As in \Cref{thm:defect-decomposition}(b):
$\int_{|z|\leq 1}|z|^p\,\nu(dz) = \infty$ for $p < \gamma$,
so $\lambda_t z \notin L^p(\nu)$ for any $\lambda_t \neq 0$.

(b) This is the key estimate.  By~\eqref{eq:psi-def},
$\psi_t^{\mathrm{def}}(z)$ is the conditional expectation of
the second-order-and-higher response.  We bound it by tracing
the composition chain~\eqref{eq:composition-chain}.

\textbf{Step 1.  Volatility shift.}  A jump of size $z$ at
time $t$ shifts $V_s \mapsto V_s + K(s,t)z$ for $s \geq t$,
hence
\begin{equation}\label{eq:sigma-shift}
  \sigma_s^{t,z} = \sigma_s\,\exp(\eta\,K(s,t)\,z),
  \qquad s \geq t.
\end{equation}
For $|z| \leq 1$ and bounded kernel: $|\sigma_s^{t,z}
- \sigma_s - \eta K(s,t)z\,\sigma_s|
\leq C\,\eta^2 K(s,t)^2 z^2\,\sigma_s$ by the Taylor
remainder of $e^x$.

\textbf{Step 2.  Stock price shift.}  The log stock price
satisfies $\log S_T = \log S_0 + \int_0^T(\mu - \frac{1}{2}
\sigma_s^2)\,ds + \int_0^T \sigma_s\,dW_s^S$.
Under the jump shift:
\begin{align}
  \log S_T^{t,z} - \log S_T^{t,0}
  &= -\tfrac{1}{2}\int_t^T
  [(\sigma_s^{t,z})^2 - \sigma_s^2]\,ds
  + \int_t^T [\sigma_s^{t,z} - \sigma_s]\,dW_s^S
  \notag\\
  &= \int_t^T \sigma_s\,(e^{\eta K(s,t)z} - 1)\,dW_s^S
  - \tfrac{1}{2}\int_t^T \sigma_s^2\,
  (e^{2\eta K(s,t)z} - 1)\,ds.
  \label{eq:logS-shift}
\end{align}
The key point: both terms are \emph{nonlinear in $z$} through
the exponential $e^{\eta K(s,t)z}$.  The first-order terms
(linear in $z$) are absorbed into $\psi^{\mathrm{lin}}$.  The
remainders are quadratic and higher in $z$, each weighted by
$K(s,t)$.

\textbf{Step 3.  Payoff shift.}  For $h \in C_b^3$,
$h(S_T^{t,z}) - h(S_T^{t,0}) - h'(S_T^{t,0})(S_T^{t,z}
- S_T^{t,0})$ is bounded by
$\frac{1}{2}\|h''\|_\infty|S_T^{t,z} - S_T^{t,0}|^2$.
The first-order term $h'(S_T)(S_T^{t,z} - S_T^{t,0})$ is
linear in $z$ (to leading order) and is absorbed into
$\psi^{\mathrm{lin}}$.

\textbf{Step 4.  Aggregate quadratic bound.}  Combining
Steps~1--3 and taking conditional expectations, the defect
satisfies, for $|z| \leq 1$:
\begin{equation}\label{eq:psi-def-bound}
  |\psi_t^{\mathrm{def}}(z)|
  \leq C(h,\eta,\sigma_0)\;\left(\int_t^T K(s,t)^2\,ds\right)
  \;z^2,
\end{equation}
where the constant absorbs $\|h'\|_\infty$, $\|h''\|_\infty$,
$\eta^2$, and moments of $\sigma$.  The factor
$\int_t^T K(s,t)^2\,ds$ is the \emph{integrated kernel weight}:
the cumulative effect of a time-$t$ jump on the volatility
path over $[t,T]$.

For $|z| > 1$: $|\psi_t^{\mathrm{def}}(z)| \leq
C(1 + |z|^p)$ by the boundedness of $h$.

\textbf{Step 5.  $\widetilde{H}_J$ norm.}
Raising~\eqref{eq:psi-def-bound} to the $p$-th power and
integrating:
\begin{align}
  \|\psi^{\mathrm{def}}\|_{\widetilde{H}_J}^p
  &\leq C^p\,\E\!\int_0^T\!\left(\int_t^T K(s,t)^2\,ds
  \right)^p\!dt\;\cdot\;\int_{|z|\leq 1}|z|^{2p}\,\nu(dz)
  \;+\;R_{\mathrm{large}}
  \notag\\
  &\leq C^p\;\|K\|_{\mathrm{agg}}^p\;
  \frac{2c_\gamma}{2p-\gamma}
  \;+\;R_{\mathrm{large}},
  \label{eq:psi-def-norm}
\end{align}
where $\|K\|_{\mathrm{agg}}^p := \int_0^T (\int_t^T
K(s,t)^2\,ds)^p\,dt$ is the \emph{aggregated kernel norm}.
The small-jump integral converges since $2p > \gamma$
(\Cref{thm:defect-decomposition}).
\end{proof}


\subsubsection{Aggregated kernel norm for fractional kernels}

\begin{lemma}[Aggregated kernel norm]\label{lem:agg-kernel}
For $K_H(s,t) = (s-t)^{H-1/2}$:
\begin{equation}\label{eq:agg-norm}
  \|K_H\|_{\mathrm{agg}}^p
  = \int_0^T\!\left(\int_t^T (s-t)^{2H-1}\,ds\right)^p dt
  = \int_0^T\!\left(\frac{(T-t)^{2H}}{2H}\right)^p dt
  = \frac{T^{2pH+1}}{(2H)^p(2pH+1)}.
\end{equation}
\end{lemma}

\begin{proof}
The inner integral is $\int_t^T(s-t)^{2H-1}\,ds =
(T-t)^{2H}/(2H)$.  Raising to the $p$-th power and
integrating: $\int_0^T(T-t)^{2pH}/(2H)^p\,dt =
T^{2pH+1}/[(2H)^p(2pH+1)]$.
\end{proof}

\begin{remark}[Comparison with the direct defect scaling]
\label{rem:comparison}
The aggregated kernel norm involves $\int_t^T K(s,t)^2\,ds$
raised to the $p$-th power, while the direct
defect (\Cref{thm:scaling}) involves $\int_0^T K(T,t)^{2p}\,dt$.
These differ:
\begin{enumerate}[label=(\roman*)]
\item The \emph{direct defect} (\Cref{sec:scaling}) applies to
  $f(V_T)$ where $V_T = \int_0^T K(T,s)\,dL_s$ is a
  \emph{terminal} Volterra integral.  The kernel weight is
  $K(T,t)^{2p}$, evaluated at the terminal time.
\item The \emph{hedging defect} (this section) applies to
  $h(S_T)$ where $S_T$ depends on the \emph{entire volatility
  path} $\{\sigma_s : s \in [0,T]\}$, which depends on
  $\{V_s : s \in [0,T]\}$.  A jump at time $t$ affects
  $V_s$ for all $s \geq t$ through $K(s,t)$, so the relevant
  weight is the \emph{integrated} kernel $\int_t^T K(s,t)^2\,ds$.
\end{enumerate}
This distinction is fundamental: path-dependent payoffs
(including any payoff on $S_T$ when the stock price is a
diffusion with stochastic volatility) are sensitive to the
kernel's integrated effect, not just its terminal value.
For the fractional kernel, $\int_t^T K_H(s,t)^2\,ds =
(T-t)^{2H}/(2H)$, while $K_H(T,t)^{2p} = (T-t)^{2p(H-1/2)}$.
Both capture the roughness, but with different exponents.
\end{remark}


\subsection{The hedging decomposition: main result}

The preceding subsections establish each step of the
hedging decomposition.  The regularity of the scenario map
$z \mapsto S_T^{t,z}$ is proved in \Cref{lem:scenario-reg}
under the exponential moment condition; the variance swap
loading is established in \Cref{lem:xi-loading}; the
integrability dichotomy and uniform $L^p$-estimates
follow from the composition chain analysis of
\Cref{prop:integrability}.  We now combine these into the
main hedging result.

\begin{theorem}[Hedging decomposition]\label{thm:hedging-main}
Under \Cref{ass:market,ass:nu,ass:kernel,ass:exp-moment},
let $h \in C_b^3(\R)$
and let $\QQ$ be any equivalent martingale measure for $S$.  Then:
\begin{equation}\label{eq:hedging-decomp}
  B_T^{-1}h(S_T)
  = \E^\QQ[B_T^{-1}h(S_T)]
  + \underbrace{\int_0^T \Delta_t\,d\tilde{S}_t
  + \int_0^T \Delta_t^\Xi\,d\tilde{\Xi}_t}_{%
  \text{continuous hedge (stock $+$ variance swap)}}
  + \underbrace{\int_0^T\!\!\int_\R
  \psi_t^{\mathrm{def}}(z)\,\Ntilde^\QQ(dt,dz)
  + R_T}_{\text{jump residual}},
\end{equation}
where:
\begin{enumerate}[label=(\alph*)]
\item $\tilde{S}_t = e^{-rt}S_t$ and $\tilde{\Xi}_t$ are the
  discounted stock and variance swap.  The hedging ratios
  $\Delta_t$ (stock delta) and $\Delta_t^\Xi$ (variance swap
  loading) are predictable and determined by the continuous
  MRT integrands $\phi^\perp, \phi^W$, the leverage $\rho$,
  and the variance swap loading $\xi^W$.
\item $\psi^{\mathrm{def}} \in \widetilde{H}_J$ is the
  $\nu$-regular defect (\Cref{prop:integrability}(b)), with
  norm bounded by~\eqref{eq:psi-def-norm}.
\item $R_T = \int_0^T\!\!\int_\R \psi_t^{\mathrm{lin}}(z)\,
  \Ntilde^\QQ(dt,dz)$ is the $\nu$-singular first-order jump
  exposure.  This is a well-defined $L^p$-martingale (despite
  $\lambda z \notin \widetilde{H}_J$), since the
  \emph{one-parameter} integral $\int \lambda_t\,dL_t$
  converges in $L^p$ by the BDG inequality for the one-parameter
  energy space $H_L$.
\end{enumerate}
\end{theorem}

\begin{proof}
Start from the three-driver MRT (\Cref{prop:MRT}).  Apply
\Cref{lem:after-delta}: stock delta-hedging absorbs
$\phi^S\,dW^{S,\QQ}$.
Apply \Cref{prop:residual-SV}: variance swap hedging absorbs
$\phi^{W,\perp}\,dW^{\perp_\star,\QQ}$ and contributes a
jump correction.  The remaining Poisson integral
$\int\!\!\int \psi_t^{\mathrm{total}}(z)\,\Ntilde^\QQ(dt,dz)$
decomposes into $\psi^{\mathrm{lin}} + \psi^{\mathrm{def}}$ by
\Cref{thm:jump-decomp}.  The norm bound on $\psi^{\mathrm{def}}$
follows from \Cref{prop:integrability}.
\end{proof}


\begin{corollary}[Hedging error scaling]\label{cor:hedge-error}
Under the fractional kernel $K_H$ with $H \in (0,1)$ and
$\gamma$-stable driver with $\gamma \in (1,2)$, the
$L^p$-norm of the $\nu$-regular defect satisfies:
\begin{equation}\label{eq:hedge-error}
  \left\|\int_0^T\!\!\int_\R \psi^{\mathrm{def}}\,
  d\Ntilde^\QQ\right\|_{L^p}
  \leq C_p\,
  \|\psi^{\mathrm{def}}\|_{\widetilde{H}_J}
  \leq C(h,\eta,p)\;\frac{T^{2H + 1/p}}{(2H)^{1+1/p}
  (2pH+1)^{1/p}}
  \;\cdot\;\left(\frac{2c_\gamma}{2p - \gamma}\right)^{1/p}
  + R^{1/p},
\end{equation}
where the exponent of $T$ is $2H + 1/p > 1/p$, confirming that
roughness ($H > 0$) amplifies the defect relative to the
white-noise limit.
\end{corollary}

\begin{proof}
Apply the BDG inequality to bound
$\|\widetilde{\delta}_J(\psi^{\mathrm{def}})\|_{L^p}$
by $C_p\|\psi^{\mathrm{def}}\|_{\widetilde{H}_J}$.
Substitute~\eqref{eq:psi-def-norm} with the aggregated kernel
norm from \Cref{lem:agg-kernel}, and take $p$-th roots.
\end{proof}

\begin{remark}[Financial interpretation]\label{rem:hedging-practical}
For a practitioner running a book with rough-jump volatility
models, the bound~\eqref{eq:hedge-error} gives the size of the
residual risk after delta- and vega-hedging.  The key inputs are
three numbers that can be estimated from data: the Hurst
parameter~$H$ (from realized volatility scaling), the stable
index~$\gamma$ (from return tail estimates), and the
horizon~$T$.  The Markov case $H = 1/2$ gives the minimum
hedging error; every rough model ($H < 1/2$) increases the
residual strictly.  The bound is computable and can be used for
reserve calculations.
\end{remark}


\subsection{The boundary of hedgeability}

\Cref{thm:hedging-main} identifies three levels of
hedgeability, each requiring progressively richer
instrument sets:

\begin{enumerate}[label=\textbf{Level \arabic*.}]
\item \textbf{Continuous instruments} (stock + variance swap).
  These absorb the full Brownian part
  ($\phi^\perp\,dW^{\perp,\QQ}$ and $\phi^W\,dW^\QQ$).
  The residual is the total Poisson integral.

\item \textbf{First-order jump instruments} (the L\'evy driver
  $L$ itself, or a sufficiently liquid proxy).  Trading $L$
  absorbs the $\nu$-singular component
  $R_T = \int\lambda_t z\,\Ntilde^\QQ(dt,dz)$, the
  first-order jump exposure.  The residual is the
  $\nu$-regular defect $\Gamma_V^h$.

\item \textbf{Nonlinear jump instruments} (options on
  jump-squared, variance-of-jumps contracts).  Since
  $\psi^{\mathrm{def}} \in \im(\widetilde{\delta}_J)$, the
  defect is in principle representable as a
  $\Ntilde$-integral with known integrand---hence hedgeable
  if instruments sensitive to $g(z) \sim z^2$ are available.
\end{enumerate}

\begin{proposition}[Vanishing conditions]\label{prop:vanishing}
The defect $\psi^{\mathrm{def}}$ vanishes:
\begin{enumerate}[label=(\alph*)]
\item if $\nu = 0$ (no jumps): $\psi^{\mathrm{def}} = 0$
  identically, recovering perfect hedging with stock +
  variance swap in pure rough volatility models
  (\cite{ElEuchRosenbaum2019,Fukasawa2021});
\item if $h$ is affine: $h(S) = a + bS$ implies
  $h(S_T^{t,z}) - h(S_T^{t,0}) = b(S_T^{t,z} - S_T^{t,0})$,
  which is linear in $z$ (to leading order through the
  exponential map), and the second-order remainder is
  absorbed into $\psi^{\mathrm{def}}$ but with
  $\|h''\|_\infty = 0$, so the bound~\eqref{eq:psi-def-bound}
  degenerates.  For $h(S) = S$ (the stock itself), the
  hedging error is zero by definition of the EMM.
\end{enumerate}
The defect is \emph{generically nonzero} for nonlinear payoffs
($h'' \neq 0$) in the presence of jumps ($\nu \neq 0$) with
rough volatility ($H \neq 1/2$): all three ingredients
contribute to~\eqref{eq:psi-def-bound}.
\end{proposition}

\begin{remark}[The multiplicative structure, revisited]
\label{rem:multiplicative}
The bound~\eqref{eq:psi-def-bound} makes the multiplicative
interaction explicit at the integrand level:
\[
  |\psi_t^{\mathrm{def}}(z)| \leq
  C \cdot \underbrace{\vphantom{\bigg|}
  \int_t^T K(s,t)^2\,ds}_{\text{kernel (roughness)}}
  \;\cdot\; \underbrace{z^2\vphantom{\bigg|}}_{\text{jump}}.
\]
Neither factor alone produces an unhedgeable residual:
without the kernel ($K = 1$, Markov case), the integrated weight
is $T - t$ and the defect is the standard L\'evy Leibniz defect.
Without jumps ($z^2 \to 0$), the defect vanishes regardless of
the kernel.  The kernel \emph{amplifies} the jump contribution
by weighting each jump according to its temporal persistence in
the volatility dynamics---recent jumps ($t$ near $T$) have small
weight $\int_t^T K^2\,ds$ for non-singular kernels, but for
rough kernels with $H < 1/2$, $K(s,t) = (s-t)^{H-1/2}$ is
singular and the weight is large near $s = t$, concentrating
the defect on short-lived but intense volatility impacts.
\end{remark}


%======================================================================
\section{Numerical Verification}\label{sec:numerics}
%======================================================================

We verify the analytical results of
\Cref{sec:pricing,sec:scaling,sec:hedging} numerically,
using the quadratic payoff $f(x) = x^2$ (where the defect
integrand is exact and deterministic,
\Cref{thm:sharpness}) and the full rough-jump stochastic
volatility model of \Cref{sec:hedging}.

Throughout, the L\'evy measure is symmetric tempered stable:
$\nu(dz) = c_\gamma\,|z|^{-(1+\gamma)}\,e^{-\lambda|z|}\,dz$
with $c_\gamma = 1$, $\lambda = 1$ (unless noted otherwise),
and the working exponent is $p = 3/2$.
Stochastic experiments use $4 \times 10^5$ paths on an
NVIDIA RTX~5090 GPU (34.2~GB VRAM) via PyTorch in
\texttt{float32} arithmetic.
Deterministic experiments evaluate the exact formulas of
\Cref{thm:scaling,thm:sharpness} directly.


\subsection{Scaling law}

\Cref{thm:scaling} predicts that the defect $L^p$-norm
scales as
$\|\Gamma_V^{x^2}\|_{L^p} \sim C \cdot T^{2H - 1 + 1/p}
\cdot (2p - \gamma)^{-1/p}$.
We verify both exponents.

\subsubsection{$T$-scaling}

For $\gamma = 1.5$ and $H \in \{0.15, 0.25, 0.35, 0.45, 0.50\}$,
we evaluate $\|\Gamma_V^{x^2}\|_{L^p}$ (via the exact
formula~\eqref{eq:gfK-exact-norm}) over
$T \in [10^{-1}, 10^2]$.  On a log-log plot, the slopes
agree with the predicted exponent $2H - 1 + 1/p$ to
machine precision (Figure~\ref{fig:scaling}a).

\subsubsection{$\gamma$-scaling}

For $H = 0.3$, $T = 1$, we sweep $\gamma \in (1.05, 2p^-)$
and compare the defect norm to the leading-order
prediction $(2p - \gamma)^{-1/p}$.
The two curves are indistinguishable up to a
multiplicative constant, confirming the blow-up rate as
$\gamma \to 2p^-$ (Figure~\ref{fig:scaling}b).
The large-jump remainder $R(f,T,p)$
in~\eqref{eq:scaling-law} is subdominant throughout the
range.

\begin{figure}[ht]
\centering
% [Include fig1_scaling_law.pdf]
\caption{Scaling law verification (\Cref{thm:scaling}).
(a)~Log-log plot of $\|\Gamma_V^{x^2}\|_{L^p}$ vs.\ $T$
for five values of~$H$; slopes match the theoretical
exponent $2H - 1 + 1/p$.
(b)~$\gamma$-dependence at $H = 0.3$, $T = 1$: the
exact defect norm (solid) vs.\ the leading-order
$(2p - \gamma)^{-1/p}$ (dashed), confirming the
blow-up rate.  Parameters: $c_\gamma = 1$,
$\lambda = 1$, $p = 3/2$.}
\label{fig:scaling}
\end{figure}


\subsection{Multiplicative interaction}

\Cref{cor:interaction} asserts that roughness and jump
intensity interact multiplicatively:
$\|\Gamma_V^f\|_{L^p} \lesssim
T^{2H - 1 + 1/p} \times (2p - \gamma)^{-1/p}$.
Figure~\ref{fig:heatmap} displays
$\|\Gamma_V^{x^2}\|_{L^p}$ over
$(H, \gamma) \in (0.02, 0.50) \times (1.05, 1.98)$
at $T = 1$.

The heatmap exhibits a $35\times$ dynamic range, with the
defect largest at small $H$ (rough) and $\gamma$ near $2p$
(heavy-tailed).  The contour lines are approximately
hyperbolic, confirming the product structure
$\log\|\Gamma\| \approx a(H) + b(\gamma)$.
We verify this quantitatively by plotting
$\log\|\Gamma(H,\gamma)\| - \log\|\Gamma(H_{\mathrm{ref}},\gamma)\|$
as a function of $\gamma$ for several values of $H$:
the curves are nearly constant, confirming
$\log$-separability (Figure~\ref{fig:multiplicative}b).

\begin{figure}[ht]
\centering
% [Include fig2_heatmap.pdf]
\caption{Kernel-weighted Leibniz defect
$\|\Gamma_V^{x^2}\|_{L^p}$ over
$(H, \gamma)$.  The dashed line marks the Markov case
$H = 1/2$.  Contour lines are hyperbolic, confirming
the multiplicative interaction of
\Cref{cor:interaction}.  Parameters: $T = 1$,
$c_\gamma = 1$, $\lambda = 1$, $p = 3/2$.}
\label{fig:heatmap}
\end{figure}


\subsection{Hedging P\&L}

We simulate $10^5$ paths of the rough-jump stochastic
volatility model~\eqref{eq:S-dynamics}--\eqref{eq:V-dynamics}
with parameters:
\[
  H = 0.3,\quad \gamma = 1.5,\quad \rho = -0.7,\quad
  \eta = 0.5,\quad \sigma_L = 0.3,\quad \sigma_0 = 0.20,\quad
  T = 1,\quad S_0 = K = 100.
\]
The L\'evy driver uses tempered stable jumps ($c_\gamma = 0.5$,
$\lambda = 2$) with a compound Poisson superposition for
$|z| \geq 0.1$ and a Gaussian correction for $|z| < 0.1$.
The Volterra process $V_t = \int_0^t K_H(t,s)\,dL_s$ is
computed by a single GPU matrix multiply
($10^5 \times 500$ increments against the $500 \times 500$
kernel matrix).

We compare three hedging strategies for an at-the-money
European call:

\begin{enumerate}[label=(\roman*)]
\item \textbf{Stock only} (delta hedge with Black--Scholes
  delta at current implied vol):
  hedging error std $= 4.52$.
\item \textbf{Stock $+$ variance swap} (regression-based
  vega hedge against realized variance):
  hedging error std $= 2.74$, a 39\% reduction.
\item \textbf{Stock $+$ variance swap $+$ jump proxy}
  (excess quadratic variation
  $\sum (\Delta\log S)^2 - \int \sigma^2\,dt$ as proxy
  for the compensated jump variation):
  hedging error std $= 2.74$, negligible further reduction.
\end{enumerate}

The transition from Level~1 to Level~2 confirms the
continuous--jump hierarchy of \Cref{thm:hedging-main}:
stock and variance swap absorb the Brownian part, reducing
the error by $39\%$.  The negligible improvement from the
jump proxy reflects an important practical point: the
Level~3 instruments in \Cref{thm:hedging-main}
(``options on jump-squared'') are structurally available
but lack liquid market analogs.  The excess quadratic
variation of stock returns is a noisy estimator of the
compensated jump quadratic variation of the volatility
process; it captures the wrong object.
A direct proxy for $\int\!\!\int z^2\,\Ntilde(dt,dz)$
in the volatility dynamics would require an instrument
sensitive to the jump variation of volatility itself,
not of the stock price.

\begin{figure}[ht]
\centering
% [Include fig3_hedging_pnl.pdf]
\caption{Hedging error distributions ($10^5$ paths).
Stock-only (left), stock $+$ variance swap (center),
stock $+$ variance swap $+$ jump proxy (right).
The variance swap absorbs the continuous volatility
exposure (39\% std reduction);
the jump proxy provides negligible further improvement,
consistent with the absence of liquid Level~3 instruments.
Parameters as in the text.}
\label{fig:hedging}
\end{figure}


\subsection{Recovery checks}

\subsubsection{Markov recovery}

As $H \to 1/2$, the kernel $K_H(T,t) = (T-t)^{H-1/2}$
converges to $K_{1/2} \equiv 1$, and the defect should
approach the standard (non-Volterra) L\'evy Leibniz
defect of~\cite{Fontes2026d}.
Figure~\ref{fig:recovery}(a) plots
$\|\Gamma_V^{x^2}\|_{L^p}$ as a function of $H$
for $\gamma = 1.5$, $T = 1$: the defect decreases
monotonically to the Markov level as $H \to 1/2$,
confirming \Cref{cor:interaction}(b).

\subsubsection{Pure rough-volatility recovery}

As $\lambda \to \infty$ (tempering parameter), the
L\'evy measure decays exponentially and $\nu \to 0$,
recovering pure rough volatility.  The defect norm
should vanish, and the rate is
$\|\Gamma_V^{x^2}\|_{L^p} \sim \lambda^{(\gamma - 2p)/p}$
(from $\nu_{2p} \sim \lambda^{\gamma - 2p}$).
Figure~\ref{fig:recovery}(b) confirms this on a log-log
plot: the slope matches the theoretical exponent
$(\gamma - 2p)/p = -1.0$ to four decimal places.

\begin{figure}[ht]
\centering
% [Include fig3_recovery.pdf]
\caption{Recovery checks.
(a)~$H \to 1/2$: defect approaches the Markov level
(dashed).
(b)~$\lambda \to \infty$ ($\nu \to 0$): defect vanishes
with slope $(\gamma - 2p)/p = -1.0$ (dashed), confirming
the pure-rough-vol recovery (\Cref{prop:vanishing}(a)).
Parameters: $T = 1$, $c_\gamma = 1$, $p = 3/2$.}
\label{fig:recovery}
\end{figure}


\subsection{$L^2$ isometry}

\Cref{thm:sharpness}(c) gives the exact identity
$\E[|\Gamma_V^{x^2}|^2]
= T^{4H-1}/(4H-1) \cdot \nu_4$
for deterministic integrands (the It\^o isometry).
We verify this by continuous-time Monte Carlo:
for each path, we sample $N \sim \mathrm{Poisson}(\Lambda T)$
exact jump times $t_j \sim \mathrm{Uniform}(0,T)$ and
magnitudes from the tempered stable distribution
(compound Poisson, $|z| \geq 0.1$), evaluate
$K(T,t_j)^2 z_j^2$ at each jump, and form
$\Gamma = \sum_j K(T,t_j)^2 z_j^2 - \nu_{2,\mathrm{cp}} \cdot
\int_0^T K(T,t)^2\,dt$
(where $\nu_{2,\mathrm{cp}} = \int_{|z| \geq 0.1}
z^2\,\nu(dz)$).
The continuous-time formulation avoids the Riemann-sum
bias that arises when $K(T,t)^4 = (T-t)^{4H-2}$ is
integrated on a time grid: for $H = 0.28$, the integrand
has exponent $4H - 2 = -0.88$, and a left-endpoint sum
on $n$ points captures only fraction
$(1 - 1/n)^{4H-1}$ of the exact integral.

Table~\ref{tab:isometry} reports the results for
$\gamma = 1.5$, $\lambda = 2$, $T = 1$, with
$4 \times 10^5$ paths per $H$ value.

\begin{table}[ht]
\centering
\begin{tabular}{c|ccc}
$H$ & Theory & MC & Ratio (MC/Theory) \\
\hline
$0.28$ & $3.896$ & $3.40 \pm 0.64$ & $0.87$ \\
$0.32$ & $1.670$ & $1.70 \pm 0.11$ & $1.02$ \\
$0.38$ & $0.899$ & $0.93 \pm 0.02$ & $1.03$ \\
$0.42$ & $0.687$ & $0.71 \pm 0.01$ & $1.03$ \\
$0.50$ & $0.467$ & $0.46 \pm 0.01$ & $0.99$ \\
\end{tabular}
\caption{$L^2$ isometry verification
(\Cref{thm:sharpness}(c)):
$\E[|\Gamma_V^{x^2}|^2]$ vs.\ the theoretical value
$T^{4H-1}/(4H-1) \cdot \nu_4$.
The MC values use continuous-time Poisson simulation
($4 \times 10^5$ paths).
At $H = 0.28$, the wide confidence interval reflects the
heavy-tailed distribution of $\Gamma^2$: a small number
of jumps near $t = T$ dominate through the singular
kernel weight $K(T,t)^4 = (T-t)^{-0.88}$.
The theory value lies within the 95\% CI for all $H$.
Parameters: $\gamma = 1.5$, $\lambda = 2$, $c_\gamma = 1$.}
\label{tab:isometry}
\end{table}

For $H \geq 0.32$, the MC/Theory ratio is within $3\%$ of
unity.  At $H = 0.28$ ($4H - 1 = 0.12$, a near-critical
singularity), the ratio is $0.87$ with a wide 95\% CI of
$\pm 0.64$---the theory value $3.90$ lies well within the
interval $[2.76, 4.05]$.  The large variance at low $H$ is
physical: the kernel weight $K(T,t)^4$ concentrates on a
neighborhood of $T$ of width $\sim \epsilon$, so the sum
$\sum K^4 z^4$ is dominated by a few jumps near $t = T$,
producing heavy-tailed $\Gamma^2$.  This is consistent with
the observation that the defect concentrates on recent jumps
(\Cref{cor:interaction}(c)).





%======================================================================

%======================================================================
\section{Discussion}\label{sec:discussion}
%======================================================================

This paper establishes the pricing and hedging theory for
Volterra--L\'evy markets---financial models where the volatility
process has both memory (through a causal kernel~$K$) and jumps
(through a L\'evy driver~$L$ with measure~$\nu$).

The central contribution is the two-parameter VRNC
(\Cref{thm:VRNC-VL}): risk-neutral pricing in these
markets requires solving a Volterra integral equation for
an effective drift rate~$\Psi$ that couples the Brownian and
jump risk premia through the kernel.  This is a genuinely new
structure---in the Markov case the constraint is algebraic and
pointwise, but for non-Markov kernels the memory of the
volatility process propagates into the pricing measure.
The constraint holds for all kernel regularities~$H$, and
determines~$\Psi$ uniquely; the remaining degrees of freedom
(the continuous--jump split and the shape of the jump pricing
kernel) are characterized in \Cref{prop:freedom}.

The structural invariance of~$H$ (\Cref{thm:bracket-inv})
ensures that the roughness parameter is a genuine feature of
the market, not a calibration artifact.  Together with the VRNC,
this means that $H$ can be estimated from historical data and
used directly in the pricing framework.  The
$\nu$-singularity persistence (\Cref{prop:VT-singular})
shows that this conclusion is robust: the kernel cannot
regularize the jump-size singularity, so the
$\nu$-singular/$\nu$-regular decomposition
(\Cref{thm:defect-decomposition}) carries over unchanged.

The hedging decomposition (\Cref{thm:hedging-main}) and
the scaling law (\Cref{thm:scaling}) give the quantitative
consequences: the structurally unhedgeable component of any
European claim is the kernel-weighted Leibniz defect, whose
$L^p$~norm scales multiplicatively in the roughness and jump
parameters.  The scaling is sharp for the variance swap payoff
(\Cref{thm:sharpness}).  For a practitioner, this means:
if you are running a book with rough-jump volatility models,
the residual risk after delta- and vega-hedging has a precise,
computable size controlled by~$H$, $\gamma$, and~$T$.

From the perspective of the operator factorization program
of~\cite{Fontes2026a,Fontes2026d},
this paper demonstrates that the operator derivative
$D = \delta^*$---defined by duality, not Gaussianity---does
genuine financial work in the non-Gaussian setting.  The
Leibniz defect of~$D$ is the non-Gaussian content; in the
Volterra--L\'evy setting, this coincides exactly with the
structurally unhedgeable risk.  This identification would
not be available through Malliavin calculus, which requires
separate constructions for the Brownian and jump components
and does not produce the kernel-weighted defect as a single
algebraic object.


\subsection{Formal verification}

The algebraic spine of this paper has been formally verified
in the Lean~4 proof assistant (zero \texttt{sorry}, zero axioms),
extending the companion formalizations
of~\cite{Fontes2026a,Fontes2026d}.  The formalization carries
exactly one stochastic analysis input as a structure field:
the L\'evy--It\^o formula for semimartingales with jumps
(\cite{Applebaum2009}, Theorem~4.4.7).  All downstream
results---the operator derivative splitting, the four-term
Volterra--L\'evy decomposition, the two-regime
$\nu$-singular/$\nu$-regular splitting, the defect
identification as the Taylor remainder of~$f$ composed with
the kernel, bracket invariance, the VRNC, the scaling law
exponents, and the hedging decomposition---are proved as
theorems from this single input combined with the operator
framework.  This architecture reflects the current state of
Mathlib, where the Cameron--Martin theorem for Gaussian measures
remains under active development (Degenne, Mathlib
PR~\#26291, 264~commits), and stochastic integration for
semimartingales with jumps does not yet exist in any Lean~4
library.  The formalization and an interactive dependency
blueprint are available at
\texttt{github.com/QuijoticResearch/volterra-levy}.


\subsection{Open problems}

\begin{enumerate}[label=(\roman*)]
\item \textbf{Measure-selection within the VRNC.}
  \Cref{thm:VRNC-VL} determines~$\Psi$ but leaves the
  split $(\theta^W, \theta^J)$ free.
  Natural selection criteria include minimal entropy,
  variance-optimal hedging, and calibration to variance swap
  or VIX option prices.  Characterizing these selections and
  their effect on the hedging error constant is the most
  pressing applied problem.

\item \textbf{Path-dependent payoffs.}
  The hedging decomposition (\Cref{thm:hedging-main})
  covers European claims $h(S_T)$.  Extension to barriers,
  lookbacks, and Asians requires applying the modified Leibniz
  rule at each monitoring date,
  generating a sequence of defect terms coupled by the
  kernel's memory structure.

\item \textbf{Purely stable drivers.}
  The exponential moment condition (\Cref{ass:exp-moment})
  fails for purely $\gamma$-stable drivers.
  Either the volatility specification must be modified, or
  the hedging bounds must be derived without exponential
  moments using $L^r$-estimates for $r < \gamma$.

\item \textbf{Optimal defect decomposition.}
  The defect $\Gamma_V^f \in \mathrm{im}(\widetilde{\delta}_J)$
  is in principle representable.  Can it be further decomposed
  into hedgeable and truly unhedgeable components if structured
  instruments (options on realized jump variation) are available?
\end{enumerate}


%======================================================================
\appendix
\section{Regularity Conditions for the Stochastic Volatility
  Model}\label{app:regularity}
%======================================================================

The hedging decomposition of \Cref{sec:hedging} requires
certain integrability conditions on the stochastic
volatility model.  We collect them here.

\begin{assumption}[Exponential moment condition]
\label{ass:exp-moment}
The Volterra--L\'evy process $V$ and the exponential
specification~\eqref{eq:vol-dynamics} satisfy:
\begin{enumerate}[label=(E\arabic*)]
\item\label{E1} \textbf{Finite exponential moments:}
  $\E[\exp(r\,|V_t|)] < \infty$ for all $r > 0$ and
  $t \in [0,T]$.
  This holds if $L$ has exponentially decaying tails
  (e.g., tempered stable) and $K$ satisfies
  \Cref{ass:kernel}\ref{K1}.
  For purely stable $L$ (power-law tails), the condition
  restricts to $\eta$ small enough that
  $\E[\exp(2\eta|V_t|)] < \infty$.
\item\label{E2} \textbf{Volatility moment condition:}
  $\sup_{t \in [0,T]}\E[\sigma_t^{4p}] < \infty$ for
  the relevant $p \in (1,\gamma)$.  Under~\ref{E1}, this
  follows from the exponential specification.
\end{enumerate}
\end{assumption}

\begin{remark}[When does \ref{E1} hold?]\label{rem:E1}
For a purely $\gamma$-stable driver $L$ (no Gaussian component,
$\sigma_L = 0$), the Laplace transform $\E[e^{r|V_t|}]$
is \emph{infinite} for all $r > 0$ and $\gamma < 2$: stable
laws have no exponential moments, and the Volterra
convolution does not create them.

This means~\ref{E1} fails for purely stable drivers.  In
this case, the exponential volatility
specification~\eqref{eq:vol-dynamics} must be modified---e.g.,
using $\sigma_t = \sigma_0(1 + \eta V_t)^+$ or a bounded
specification $\sigma_t = \sigma_0\exp(\eta V_t
\mathbf{1}_{|V_t| \leq M})$.  Alternatively, the
hedging bounds of \Cref{thm:hedging-main} hold
with the constant $C(h,\eta,p)$ depending on the
volatility moments, which must be verified case by case.

For tempered stable drivers (where
$\nu(dz) = c|z|^{-1-\gamma}e^{-\lambda|z|}\,dz$ for some
$\lambda > 0$), all moments are finite and~\ref{E1}
holds for $r < \lambda$.
\end{remark}

\begin{proposition}[The bound constants are finite]
\label{prop:constants-finite}
Under \Cref{ass:exp-moment}, the constant
$C(h,\eta,p)$ in the hedging error bound~\eqref{eq:psi-def-bound}
is finite.  Specifically:
\[
  C(h,\eta,p) \leq C_0\,\|h''\|_\infty\,\eta^2\,
  \bigl(\sup_{t \in [0,T]}\E^\QQ[\sigma_t^{4p}]\bigr)^{1/(2p)},
\]
where $C_0$ depends only on $p$ and $T$.
\end{proposition}

\begin{proof}
In the bound $|\psi_t^{\mathrm{def}}(z)| \leq
C\,(\int_t^T K(s,t)^2\,ds)\,z^2$
(\Cref{prop:integrability}, Step~4),
the constant $C$ arises from the conditional expectation
of the composition chain: the volatility shift
$|\sigma_s^{t,z} - \sigma_s| \leq C\eta|K(s,t)||z|\sigma_s$,
the stock price shift
$|S_T^{t,z} - S_T^{t,0}| \leq C\,S_T\,
\int_t^T|\sigma_s^{t,z} - \sigma_s|\,|dW_s^S|$, and
the payoff shift $|h(S_T^{t,z}) - h(S_T^{t,0})
- h'(S_T^{t,0})(S_T^{t,z} - S_T^{t,0})|
\leq \frac{1}{2}\|h''\|_\infty\,|S_T^{t,z} - S_T^{t,0}|^2$.

Taking $L^p$-norms and applying H\"older's inequality,
the volatility and stock moments enter through
$\E[\sigma_s^{4p}]^{1/(2p)}$ (from the Burkholder--Davis--Gundy
inequality applied to the stock-shift stochastic integral).
Under~\ref{E2}, this is bounded.
\end{proof}


%======================================================================
% Acknowledgements
%======================================================================

\subsection*{Acknowledgements}

AI tools were used during drafting and mathematical exploration;
the author takes full responsibility for all content and errors.


%======================================================================
% References
%======================================================================

\begin{thebibliography}{9}

\bibitem{Applebaum2009}
D.~Applebaum, \emph{L\'evy Processes and Stochastic Calculus},
2nd ed., Cambridge University Press, 2009.

\bibitem{ContTankov2004}
R.~Cont and P.~Tankov, \emph{Financial Modelling with Jump
Processes}, Chapman \& Hall/CRC, 2004.

\bibitem{Dirksen2015}
S.~Dirksen, \emph{It\^o isomorphisms for $L^p$-valued Poisson
stochastic integrals}, Ann.\ Probab.\ 43 (2015), 2793--2825.

\bibitem{ElEuchRosenbaum2019}
O.~El~Euch and M.~Rosenbaum, \emph{Perfect hedging in rough
Heston models}, Ann.\ Appl.\ Probab.\ 29 (2019), 3223--3268.

\bibitem{Fontes2026a}
R.~Fontes, \emph{The operator derivative in continuous stochastic
calculus: a Hilbert energy space framework}, Quijotic Research,
2026.

\bibitem{Fontes2026d}
R.~Fontes, \emph{Operator factorization beyond Hilbert spaces:
representability obstructions and Leibniz defects for stable
L\'evy processes}, Quijotic Research, 2026.

\bibitem{Fukasawa2021}
M.~Fukasawa and J.~Gatheral, \emph{Hedging under rough
volatility}, in Rough Volatility, SIAM, 2021.

\bibitem{GripenbergLondenStaffans1990}
G.~Gripenberg, S.-O.~Londen, and O.~Staffans,
\emph{Volterra Integral and Functional Equations},
Cambridge University Press, 1990.

\end{thebibliography}


\end{document}